\documentclass[11pt,a4paper]{article}

\usepackage[T1]{fontenc}
\usepackage[utf8]{inputenc}
\usepackage{mathpazo}
\usepackage[scaled=0.85]{helvet}
\renewcommand{\ttdefault}{cmtt}
\usepackage{microtype}
\usepackage[margin=1.05in,top=1.1in,bottom=1.15in]{geometry}
\usepackage{amsmath,amssymb,amsthm,mathtools}
\usepackage{booktabs}
\usepackage{enumitem}
\usepackage{array}
\usepackage{longtable}
\usepackage{xcolor}
\usepackage{titlesec}
\usepackage[hidelinks,colorlinks=true,linkcolor=linknavy,citecolor=linknavy,urlcolor=linknavy]{hyperref}

\definecolor{linknavy}{HTML}{1F3A63}
\definecolor{rulegray}{HTML}{888888}
\definecolor{boxbg}{HTML}{F4F2ED}

% ---------------------------------------------------------------- theorem envs
\theoremstyle{plain}
\newtheorem{theorem}{Theorem}[section]
\newtheorem{proposition}[theorem]{Proposition}
\newtheorem{lemma}[theorem]{Lemma}
\newtheorem{corollary}[theorem]{Corollary}

\theoremstyle{definition}
\newtheorem{definition}[theorem]{Definition}
\newtheorem{guarantee}[theorem]{Architectural guarantee}
\newtheorem{observation}[theorem]{Numerical observation}

\theoremstyle{remark}
\newtheorem{remark}[theorem]{Remark}
\newtheorem{trap}[theorem]{Trap}

\renewcommand{\qedsymbol}{$\blacksquare$}

% ---------------------------------------------------------------- section look
\titleformat{\section}{\normalfont\Large\bfseries}{\thesection}{0.8em}{}
\titleformat{\subsection}{\normalfont\large\bfseries}{\thesubsection}{0.7em}{}
\titleformat{\subsubsection}{\normalfont\bfseries}{\thesubsubsection}{0.6em}{}
\titlespacing*{\section}{0pt}{2.4ex plus 1ex minus .2ex}{1.4ex plus .2ex}

% ---------------------------------------------------------------- macros
\newcommand{\R}{\mathbb{R}}
\newcommand{\C}{\mathbb{C}}
\newcommand{\Z}{\mathbb{Z}}
\newcommand{\T}{\mathbb{T}}
\newcommand{\iu}{\mathrm{i}}
\newcommand{\eu}{\mathrm{e}}
\newcommand{\dd}{\,\mathrm{d}}
\newcommand{\lap}{\Delta}
\newcommand{\psih}{\hat{\psi}}
\newcommand{\dt}{\Delta t}
\newcommand{\half}{\tfrac{1}{2}}
\newcommand{\kwrap}{k_{\mathrm{wrap}}}
\newcommand{\ktrain}{k_{\mathrm{train}}}
\newcommand{\knyq}{k_{\mathrm{Nyq}}}
\newcommand{\epssplit}{\varepsilon_{\mathrm{split}}}
\newcommand{\Kop}{\mathcal{K}}
\newcommand{\Nop}{\mathcal{N}}
\newcommand{\Flow}{\Phi}
\newcommand{\Ham}{\mathcal{H}}
\newcommand{\Mass}{\mathcal{M}}
\newcommand{\Mom}{\mathcal{P}}
\newcommand{\code}[1]{\texttt{\small #1}}
\newcommand{\src}[1]{{\color{linknavy}\texttt{\footnotesize #1}}}
\DeclareMathOperator{\Imag}{Im}
\DeclareMathOperator{\Real}{Re}
\DeclareMathOperator*{\argmin}{arg\,min}
\DeclareMathOperator{\Id}{Id}
\DeclareMathOperator{\wrap}{wrap}

\newcommand{\sourcenote}[1]{%
  \par\vspace{0.35em}\noindent
  {\footnotesize\color{rulegray}\rule{\linewidth}{0.4pt}\\[0.25em]
  \textsc{in code:} \src{#1}}\par\vspace{0.5em}}

\setlength{\parskip}{0.35em}
\setlength{\parindent}{0pt}
\begin{document}

\begin{titlepage}
\centering
\vspace*{2.2cm}

{\Large\scshape The Mathematics of SPNO}\\[0.9em]

{\LARGE\bfseries Spectral Identifiability and\\[0.25em]
Structure-Preserving Neural Operators\\[0.35em]
for the Parametric Nonlinear Schr\"odinger Equation}\\[2.2em]

{\large A derivation-level reference for the ten-phase study}\\[0.6em]
{\normalsize organised by subject, with a phase map in \S\ref{sec:phasemap}}\\[3.2em]

\begin{minipage}{0.82\textwidth}\small
Every claim in this document is labelled by \emph{kind}, following the house rule of the
project: a \textbf{Theorem} or \textbf{Proposition} is a mathematical statement with a
proof; an \textbf{Architectural guarantee} holds for every parameter value $\theta$,
untrained weights included, and is asserted in the test suite at random initialisation;
a \textbf{Numerical observation} is something that was measured and could have come out
otherwise. The distinction is load-bearing. ``Conserves energy'' is never written where
the truth is ``phase-only substeps preserve the discrete $L^2$ norm by construction, and
Hamiltonian agreement is empirical.''

\vspace{0.9em}
Numerical constants quoted in the text are the measured values from configuration hash
\code{bd4e108527} ($N=64$, $\dt=0.01$, $M=32$ substeps) unless stated otherwise, and are
collected in \S\ref{sec:constants}.
\end{minipage}

\vfill
{\small Generated from the \code{spno} source tree \textemdash{} \today}
\end{titlepage}

\pagenumbering{roman}
\tableofcontents
\newpage
\pagenumbering{arabic}

\section*{How to read this}
\addcontentsline{toc}{section}{How to read this}

The study asks one question: \emph{does encoding physical structure directly in a neural
operator's architecture beat an unconstrained baseline, or softer alternatives?} Four
ways of using physics are kept rigorously distinct, because they have different
mathematical status:

\begin{center}
\begin{tabular}{@{}p{0.26\linewidth}p{0.68\linewidth}@{}}
\toprule
\textbf{nothing} & Model A, a plain FNO. No invariant is preserved; \S\ref{sec:fno}.\\
\textbf{constraint projection} & Model B, FNO composed with a metric projection onto the
mass sphere. One scalar invariant, exactly; \S\ref{sec:projection}.\\
\textbf{soft physics loss} & Model P (PINO), a Crank--Nicolson residual penalty. Nothing
is guaranteed --- and the residual itself smuggles in mass conservation;
\S\ref{sec:pino}.\\
\textbf{architecture} & Models C1/C2/C3, learned split-step operators. Mass,
reversibility, symplecticity and $U(1)$ equivariance hold by construction;
\S\ref{sec:learnedsplit}.\\
\bottomrule
\end{tabular}
\end{center}

Part~I builds the continuum and discrete theory the whole project rests on. Part~II
derives what each model class guarantees and, equally important, what it does not.
Part~III is the identifiability theory --- the thesis spine --- which is where the only
genuinely \emph{negative} theorem lives. Part~IV covers the baselines and the two
misspecification dials that turn ``structure beats FNO'' from a tautology into a
question. Part~V collects measurement machinery and the phase map.

Sections are cross-referenced to the source file that implements them, in a rule like
this one:

\sourcenote{src/spno/equations/nls.py}

\newpage
\part*{Part I \quad Foundations}
\addcontentsline{toc}{part}{Part I \; Foundations}
\section{The equation, its Hamiltonian structure, and its invariants}
\label{sec:equation}

\subsection{Statement and sign conventions}

The object of study is the \emph{parametric} nonlinear Schr\"odinger equation on the
periodic box $\T^d = [0,2\pi)^d$ (in practice $d=1$):
\begin{equation}
\label{eq:nls}
\boxed{\;\iu\,\psi_t \;+\; \alpha\,\lap\psi \;+\; \beta\,|\psi|^2\psi \;-\; V(x)\,\psi \;=\; 0\;}
\qquad \psi(\cdot,t):\T^d\to\C,\quad V:\T^d\to\R .
\end{equation}
Equivalently, solved for the time derivative,
\begin{equation}
\label{eq:nls-evol}
\psi_t \;=\; \iu\bigl(\alpha\lap\psi + \beta|\psi|^2\psi - V\psi\bigr).
\end{equation}
The parameters $(\alpha,\beta)$ are drawn per sample; $\alpha>0$ throughout, and
$\beta$ takes either sign (focusing $\beta>0$, defocusing $\beta<0$). The learning
target is the one-step solution operator
\[
\Flow_{\dt}:\;(\psi_n,\,V,\,\alpha,\,\beta)\;\longmapsto\;\psi_{n+1}\;=\;\psi(\cdot,t_n+\dt).
\]

Three sign conventions follow from \eqref{eq:nls-evol} and are used without further
comment. They are stated here because a single flipped sign invalidates every
downstream number, and each has been pinned by a test rather than by a comment.

\begin{enumerate}[leftmargin=1.4em,itemsep=0.25em]
\item \textbf{Kinetic Fourier multiplier} $\eu^{-\iu\alpha|k|^2\tau}$.
\item \textbf{Local physical-space phase} $\eu^{+\iu(\beta|\psi|^2 - V)\tau}$.
\item \textbf{Plane-wave dispersion} $\omega = \alpha|k|^2 - \beta|A|^2 + V_0$.
\end{enumerate}

\subsection{Hamiltonian structure}

\begin{definition}[Energy]
\label{def:hamiltonian}
\[
\Ham[\psi] \;=\; \int_{\T^d}\Bigl(\alpha\,|\nabla\psi|^2 \;+\; V\,|\psi|^2 \;-\; \tfrac{\beta}{2}\,|\psi|^4\Bigr)\dd x .
\]
\end{definition}

Treating $\psi$ and $\bar\psi$ as independent, \eqref{eq:nls} is the Hamiltonian system
\begin{equation}
\label{eq:hamform}
\iu\,\psi_t \;=\; \frac{\delta\Ham}{\delta\bar\psi}.
\end{equation}

\begin{proposition}
\label{prop:hamform}
Definition~\ref{def:hamiltonian} generates \eqref{eq:nls} through \eqref{eq:hamform}.
\end{proposition}

\begin{proof}
Write $\Ham = \int \alpha\,\nabla\psi\cdot\nabla\bar\psi + V\psi\bar\psi - \tfrac{\beta}{2}(\psi\bar\psi)^2$.
Varying in $\bar\psi$ and integrating the gradient term by parts (no boundary term, the
domain is periodic),
\[
\frac{\delta\Ham}{\delta\bar\psi} \;=\; -\alpha\lap\psi \;+\; V\psi \;-\; \beta|\psi|^2\psi .
\]
Substituting into \eqref{eq:hamform} gives $\iu\psi_t + \alpha\lap\psi + \beta|\psi|^2\psi - V\psi = 0$,
which is \eqref{eq:nls}.
\end{proof}

\begin{remark}
The sign of the quartic term in $\Ham$ is fixed by this computation and is easy to get
wrong: $-\tfrac{\beta}{2}|\psi|^4$, \emph{not} $+$. The implementation writes the energy
density as $\alpha|\nabla\psi|^2 + V\rho - \tfrac{\beta}{2}\rho^2$ with $\rho=|\psi|^2$.
\end{remark}

\subsection{The two conserved quantities}

\begin{definition}[Mass and momentum]
\[
\Mass[\psi] \;=\; \int_{\T^d}|\psi|^2\dd x,
\qquad
\Mom[\psi] \;=\; \Imag\int_{\T^d}\bar\psi\,\nabla\psi\dd x .
\]
\end{definition}

\begin{theorem}[Mass conservation, via Noether]
\label{thm:mass}
$\Ham$ is invariant under the global gauge rotation $\psi\mapsto \eu^{\iu c}\psi$,
$c\in\R$. Consequently $\Mass$ is conserved along solutions of \eqref{eq:nls}, for every
$V$, $\alpha$ and $\beta$.
\end{theorem}

\begin{proof}
Invariance is immediate: every term of $\Ham$ involves $\psi$ paired with $\bar\psi$, so
the phases cancel. For conservation, compute directly:
\[
\frac{\dd}{\dd t}\Mass
= \int 2\Real\bigl(\bar\psi\,\psi_t\bigr)
= \int 2\Real\Bigl(\iu\,\bar\psi\bigl(\alpha\lap\psi + \beta|\psi|^2\psi - V\psi\bigr)\Bigr).
\]
The terms $\beta|\psi|^4$ and $-V|\psi|^2$ are real, so $\iu$ times them is purely
imaginary and contributes nothing. For the remaining term, integrate by parts:
$\int \bar\psi\lap\psi = -\int|\nabla\psi|^2$, which is real, so
$\Real(\iu\alpha\cdot(\text{real}))=0$. Hence $\dd\Mass/\dd t = 0$.
\end{proof}

\begin{theorem}[Momentum conservation, only when $V\equiv 0$]
\label{thm:momentum}
If $V\equiv 0$ then \eqref{eq:nls} is invariant under translations and $\Mom$ is
conserved. For non-constant $V$ it is not.
\end{theorem}

\begin{proof}
With $V\equiv 0$ the equation has no explicit $x$-dependence, so
$\psi(\cdot,t)\mapsto\psi(\cdot-a,t)$ maps solutions to solutions; Noether's theorem for
the translation group gives conservation of $\Mom$. Concretely, in $d=1$,
\[
\frac{\dd\Mom}{\dd t}
= \Imag\int\bigl(\bar\psi_t\psi_x + \bar\psi\psi_{xt}\bigr)
= -\int V_x\,|\psi|^2 ,
\]
after substituting \eqref{eq:nls-evol} and integrating by parts; the kinetic and
nonlinear terms produce exact derivatives, which integrate to zero on $\T$. The right-hand
side vanishes for every $\psi$ precisely when $V_x\equiv 0$.
\end{proof}

\begin{proposition}[The Strang step conserves momentum too, when $V\equiv0$]
\label{prop:discretemomentum}
On the periodic grid with spectral differentiation, $\Flow_{\dt}$ preserves the discrete
momentum $\Mom_h[\psi] = \mu\,\Imag\sum_j\overline{\psi_j}(\partial_x^{\mathrm{spec}}\psi)_j$
exactly when $V\equiv0$.
\end{proposition}

\begin{proof}
In Fourier, $\Mom_h = \tfrac{\mu}{N}\sum_k k\,|\widehat\psi_k|^2$, which the kinetic
substep leaves invariant because it changes only the \emph{phases} of $\widehat\psi_k$.
For the local substep, write $\psi\mapsto \eu^{\iu\vartheta}\psi$ with
$\vartheta = \beta\rho\,\dt$ real. Then
\[
\overline{\eu^{\iu\vartheta}\psi}\;\partial_x\bigl(\eu^{\iu\vartheta}\psi\bigr)
= \bar\psi\,\partial_x\psi \;+\; \iu\,\vartheta_x\,|\psi|^2 ,
\]
so the momentum picks up
$\int \vartheta_x\rho = \beta\dt\int\rho_x\rho = \tfrac{\beta\dt}{2}\int(\rho^2)_x = 0$,
the integral of a derivative over a periodic domain --- which vanishes discretely as well,
since the $k=0$ coefficient of a spectral derivative is exactly zero.
\end{proof}

\begin{observation}[Verified]
Over 200 steps at $\dt=0.01$ in \code{float64} on a random band-limited field:
relative momentum drift $6.7\times10^{-10}$ with $V\equiv0$ (roundoff, accumulating
linearly from $2.6\times10^{-13}$ after one step) versus $6.4\times10^{-2}$ with
$V=0.4\cos 2x$. Eight orders of magnitude of separation.
\end{observation}

\begin{remark}[Why this matters for the experiments]
The shift arm \code{G3-potential-zero} sets $V\equiv 0$, which restores translation
invariance and therefore adds a \emph{third} exactly conserved quantity to the problem
that is absent in the training distribution. That arm is consequently not just ``an
easier potential'': it is a structurally different problem, and it is the one arm where a
model could be right for a reason none of the training data exercised.
\end{remark}

\subsection{Plane waves and the dispersion relation}

\begin{proposition}[Plane wave]
\label{prop:planewave}
For \emph{constant} potential $V\equiv V_0$, the field
\[
\psi(x,t) \;=\; A\,\eu^{\iu(k\cdot x - \omega t)},\qquad A>0,\ k\in\Z^d,
\]
solves \eqref{eq:nls} exactly if and only if
\begin{equation}
\label{eq:dispersion}
\omega \;=\; \alpha|k|^2 \;-\; \beta A^2 \;+\; V_0 .
\end{equation}
\end{proposition}

\begin{proof}
Substituting: $\psi_t = -\iu\omega\psi$, $\lap\psi = -|k|^2\psi$, $|\psi|^2 = A^2$
(constant in $x$ and $t$). Then \eqref{eq:nls} reads
$\omega\psi - \alpha|k|^2\psi + \beta A^2\psi - V_0\psi = 0$, and dividing by $\psi\neq0$
gives \eqref{eq:dispersion}. The ``only if'' direction is the same computation read
backwards.
\end{proof}

Equation \eqref{eq:dispersion} is the ground truth every trained model --- black-box
models included --- is probed against in Part~III. Note the three separate dependencies
it encodes: quadratic in $k$, linear in $\alpha$, and linear in the \emph{amplitude}
through $\beta A^2$. The last one is why the probe amplitude cannot be chosen freely
(\S\ref{sec:probeamp}).

\begin{proposition}[Plane-wave mass]
\label{prop:pwmass}
$\Mass[A\eu^{\iu k\cdot x}] = A^2\prod_j L_j$, which for $\T^d=[0,2\pi)^d$ is
$A^2(2\pi)^d$.
\end{proposition}

\subsection{The discrete setting}
\label{sec:discrete}

All computation happens on a uniform, endpoint-free grid of $N$ points per dimension,
with spacing $h_j = L_j/N_j$ and cell volume $\mu = \prod_j h_j$. The discrete Fourier
transform uses the \code{norm="backward"} convention: the forward transform carries no
factor and the inverse carries $1/N$.

\begin{definition}[Discrete objects]
For a grid field $\psi\in\C^{N}$:
\[
\Mass_h[\psi] = \mu\sum_{j}|\psi_j|^2,\qquad
\widehat{\psi}_k = \sum_j \psi_j \eu^{-\iu k x_j},\qquad
\lap_h \psi = \mathcal{F}^{-1}\bigl(-|k|^2\widehat\psi\bigr),
\]
with wave vectors $k = 2\pi\,\mathrm{fftfreq}(N,h)$, so that on $[0,2\pi)$ they are exactly
the integers $\{0,1,\dots,N/2-1,-N/2,\dots,-1\}$.
\end{definition}

\begin{lemma}[Discrete Parseval]
\label{lem:parseval}
$\sum_j|\psi_j|^2 = \tfrac{1}{N}\sum_k|\widehat\psi_k|^2$. Consequently, multiplying
$\widehat\psi$ pointwise by any function of modulus one preserves $\Mass_h$.
\end{lemma}

\begin{proof}
Standard for the unnormalised DFT: $\sum_k|\widehat\psi_k|^2
= \sum_k\sum_{j,l}\psi_j\bar\psi_l\eu^{-\iu k(x_j-x_l)} = N\sum_j|\psi_j|^2$, using
$\sum_k \eu^{-\iu k(x_j-x_l)} = N\delta_{jl}$. The second claim follows because
$|m_k\widehat\psi_k| = |\widehat\psi_k|$ when $|m_k|=1$.
\end{proof}

Lemma~\ref{lem:parseval} is the entire reason the kinetic substep of every operator in
this project conserves mass exactly, and it is why the mass guarantee is
\emph{architectural} rather than merely accurate.

\begin{trap}[Wave vectors must be built in \texttt{float64} and narrowed, not built
natively]
\label{trap:k2}
At $N=64$, evaluating $2\pi\,\mathrm{fftfreq}$ in \code{float64} lands $2.3\times10^{-13}$
from the exact integers and narrows to them exactly; evaluating the same expression in
\code{float32} lands $1.2\times10^{-4}$ away, and that error \emph{survives} a later
widening to \code{float64}. Nine orders of magnitude above the $10^{-13}$ bounds the
structural tests assert. Every $|k|^2$ buffer in the codebase is therefore constructed in
double and then narrowed --- including on rebinding to a new grid.
\end{trap}

\begin{trap}[A \texttt{float32} parameter poisons a \texttt{float64} field]
\label{trap:param}
\code{torch.tensor([0.9])} is \code{float32}. Feeding it as $\alpha$ to a
\code{complex128} field costs $\approx 2.6\times10^{-8}$ in $\alpha$, which moves the
one-step phase at $k=30$ by $\approx5\times10^{-7}$ --- invalidating every
$10^{-9}$ assertion in the project. \code{batch\_parameter} rejects the downcast rather
than promoting silently.
\end{trap}

\sourcenote{src/spno/equations/nls.py, src/spno/domain.py}

\newpage
\section{Operator splitting}
\label{sec:splitting}

\subsection{The two sub-flows}

Write \eqref{eq:nls-evol} as $\psi_t = (\mathcal{A} + \mathcal{B})\psi$ with
\[
\mathcal{A}\psi \;=\; \iu\,\alpha\lap\psi
\qquad\text{(kinetic, linear, diagonal in Fourier)},
\]
\[
\mathcal{B}\psi \;=\; \iu\bigl(\beta|\psi|^2 - V\bigr)\psi
\qquad\text{(local, nonlinear, diagonal in physical space)}.
\]
Each generator has a \emph{closed-form exact flow}. That is the whole point of the
splitting: neither piece is approximated.

\begin{lemma}[Exact kinetic flow]
\label{lem:kinetic}
The solution of $\psi_t = \iu\alpha\lap\psi$ is
$\Kop_\tau\psi = \mathcal{F}^{-1}\!\left(\eu^{-\iu\alpha|k|^2\tau}\,\widehat\psi\right)$.
\end{lemma}

\begin{proof}
In Fourier, $\lap$ acts as multiplication by $-|k|^2$, so
$\partial_t\widehat\psi_k = -\iu\alpha|k|^2\widehat\psi_k$, a scalar linear ODE per mode
with solution $\widehat\psi_k(\tau)=\eu^{-\iu\alpha|k|^2\tau}\widehat\psi_k(0)$. On the
grid this is exact for every mode the grid represents; there is no spatial discretisation
error.
\end{proof}

\begin{lemma}[Exact local flow; the density is frozen]
\label{lem:local}
Along $\psi_t = \iu(\beta|\psi|^2-V)\psi$, the density $\rho = |\psi|^2$ is constant in
time \emph{pointwise}. Consequently the flow is
$\Nop_\tau\psi = \eu^{\iu(\beta\rho - V)\tau}\,\psi$, with $\rho$ evaluated at the initial
state.
\end{lemma}

\begin{proof}
$\partial_t\rho = 2\Real(\bar\psi\psi_t) = 2\Real\bigl(\iu(\beta\rho-V)|\psi|^2\bigr) = 0$,
because $\beta\rho - V$ is real and $\iu\times(\text{real})$ is purely imaginary. With
$\rho$ frozen, the equation becomes $\psi_t = \iu c(x)\psi$ with $c$ independent of $t$,
whose solution is the stated multiplier.
\end{proof}

Lemma~\ref{lem:local} is the single most important structural fact in the project. It is
what makes the local substep exactly invertible, exactly mass-preserving, and exactly
Hamiltonian --- and it is what breaks in Model C3, where the phase is allowed to read
$\Real\psi$ and $\Imag\psi$ separately (\S\ref{sec:c3}).

\subsection{The Strang step}

\begin{definition}[Symmetric / Strang splitting]
\label{def:strang}
\[
\boxed{\;\Flow_{\dt} \;=\; \Kop_{\dt/2}\circ\Nop_{\dt}\circ\Kop_{\dt/2}\;}
\]
\end{definition}

The palindromic ordering is not cosmetic. It buys three separate properties --- second
order, time-symmetry, and (with symplectic substeps) a symmetric symplectic integrator ---
none of which the asymmetric Lie--Trotter step $\Kop_{\dt}\Nop_{\dt}$ has.

\begin{theorem}[Exact mass conservation]
\label{thm:strangmass}
$\Mass_h[\Flow_{\dt}\psi] = \Mass_h[\psi]$ exactly, in exact arithmetic, for every
$\dt,\alpha,\beta,V$.
\end{theorem}

\begin{proof}
$\Nop$ multiplies each grid value by $\eu^{\iu(\beta\rho-V)\dt}$, a complex number of
modulus one (the exponent is real by Lemma~\ref{lem:local}), so $|\psi_j|$ is unchanged
at every $j$ and hence $\sum_j|\psi_j|^2$ is unchanged. $\Kop$ multiplies each Fourier
coefficient by $\eu^{-\iu\alpha|k|^2\tau}$, again of modulus one, so by
Lemma~\ref{lem:parseval} $\sum_j|\psi_j|^2$ is again unchanged. A composition of
norm-preserving maps preserves the norm.
\end{proof}

\begin{theorem}[Exact time-reversibility]
\label{thm:strangrev}
$\Flow_{-\dt}\circ\Flow_{\dt} = \Id$ exactly.
\end{theorem}

\begin{proof}
Expand, writing $h=\dt/2$:
\[
\Flow_{-\dt}\Flow_{\dt}
= \Kop_{-h}\Nop_{-\dt}\Kop_{-h}\;\Kop_{h}\Nop_{\dt}\Kop_{h}.
\]
The inner pair cancels: $\Kop_{-h}\Kop_{h}=\Id$, since the Fourier multipliers
$\eu^{+\iu\alpha|k|^2h}$ and $\eu^{-\iu\alpha|k|^2h}$ multiply to $1$ mode by mode. It
remains to show $\Nop_{-\dt}\Nop_{\dt}=\Id$. Let $w$ be the state entering $\Nop_{\dt}$
and put $\rho_w = |w|^2$. Then $v := \Nop_{\dt}w = \eu^{\iu(\beta\rho_w-V)\dt}w$, and by
Lemma~\ref{lem:local} $|v|^2 = \rho_w$ --- the backward step reads the \emph{same}
density field. Hence
$\Nop_{-\dt}v = \eu^{-\iu(\beta\rho_w-V)\dt}v = w$.
The outer pair $\Kop_{-h}\cdots\Kop_{h}$ then cancels likewise.
\end{proof}

\begin{theorem}[Symplecticity]
\label{thm:strangsympl}
$\Flow_{\dt}$ is a symplectic map on the phase space $(\psi,\bar\psi)$ with the standard
structure induced by \eqref{eq:hamform}.
\end{theorem}

\begin{proof}
By Lemma~\ref{lem:kinetic}, $\Kop_\tau$ is the exact time-$\tau$ flow of the Hamiltonian
$\Ham_K = \alpha\int|\nabla\psi|^2$; by Lemma~\ref{lem:local}, $\Nop_\tau$ is the exact
time-$\tau$ flow of $\Ham_N = \int\bigl(V\rho - \tfrac{\beta}{2}\rho^2\bigr)$. (Both
identifications are verified in the general learned setting in
Theorem~\ref{thm:learnedmain}(3), of which this is the special case $\kappa=-\alpha|k|^2$,
$\nu = \beta\rho - V$.) The exact flow of a Hamiltonian is symplectic, and a composition
of symplectic maps is symplectic.
\end{proof}

\begin{theorem}[Second order]
\label{thm:strangorder}
$\Flow_{\dt}$ is consistent of order $2$: its local error is $O(\dt^3)$.
\end{theorem}

\begin{proof}
Two routes; both are worth having.

\emph{(i) Symmetry.} By Theorem~\ref{thm:strangrev} the method satisfies
$\Flow_{-\dt}\circ\Flow_{\dt}=\Id$. Write the local error as
$\Flow_{\dt} = \varphi_{\dt} + \sum_{p\ge p_0} \dt^{\,p+1} E_p$ where $\varphi$ is the
exact flow. Time-symmetry forces the expansion of $\Flow$ in $\dt$ to contain only odd
powers beyond the exact flow, so the leading error exponent $p_0+1$ is odd, i.e.\ the
order $p_0$ is even. Since the method is consistent ($p_0\ge1$), $p_0\ge2$.

\emph{(ii) Baker--Campbell--Hausdorff.} For generators $\mathcal{A},\mathcal{B}$,
\[
\eu^{\frac{h}{2}\mathcal{A}}\eu^{h\mathcal{B}}\eu^{\frac{h}{2}\mathcal{A}}
\;=\;
\exp\Bigl(h(\mathcal{A}+\mathcal{B}) \;+\; h^3\,\mathcal{C}_2 \;+\; O(h^5)\Bigr),
\]
\begin{equation}
\label{eq:bch}
\mathcal{C}_2 \;=\; -\tfrac{1}{24}\bigl[\mathcal{A},[\mathcal{A},\mathcal{B}]\bigr]
\;+\;\tfrac{1}{12}\bigl[\mathcal{B},[\mathcal{B},\mathcal{A}]\bigr].
\end{equation}
The $h^2$ term vanishes identically by the palindromic structure --- this is exactly why
the asymmetric Lie--Trotter step is only first order --- leaving a local error $O(h^3)$.
\end{proof}

\begin{corollary}[Modified Hamiltonian and bounded energy error]
\label{cor:modham}
$\Flow_{\dt}$ is the exact flow of a modified Hamiltonian
$\widetilde\Ham = \Ham + \dt^2\Ham_2 + O(\dt^4)$, where $\dt^2\Ham_2$ is built from the
double commutators in \eqref{eq:bch}. Consequently the error in $\Ham$ along a numerical
trajectory is $O(\dt^2)$ and \textbf{bounded}, oscillating about a level rather than
growing with the horizon.
\end{corollary}

This is the falsifiable prediction that the whole rollout-diagnostics apparatus
(\S\ref{sec:drift}) exists to test. It is what separates a structure-preserving
integrator from a merely accurate one: an accurate non-symplectic method has nothing
preventing \emph{secular} energy drift.

\begin{observation}[The reference solver behaves as predicted]
\label{obs:phase0}
Over $1000$ steps at $\dt=0.01$ in \code{float64}: maximum relative mass drift
$3.82\times10^{-13}$; maximum relative energy drift $7.92\times10^{-5}$, with the
late-window maximum $0.83\times$ the early-window maximum --- classified \emph{bounded}.
Measured Strang convergence order, finest ratio: $2.0465$.
\end{observation}

\begin{trap}[Pre-asymptotic convergence ratios are meaningless]
The measured order ratios across substeps $(4,8,16,32,64)$ are
$(1.69,\,4.24,\,2.22,\,2.05)$. The $4.24$ is a cancellation artefact in the
pre-asymptotic regime, not a fourth-order method. Only the finest ratio is quoted, and
convergence claims are always made by \emph{measured scaling}
($\log_2(e_{i+1}/e_i)$ compared against the expected order) rather than by an absolute
threshold.
\end{trap}

\subsection{The substepped reference, and why targets cannot be single Strang steps}
\label{sec:substepped}

\begin{definition}[Substepped reference]
\[
\Flow^{\mathrm{ref}}_{\dt} \;=\; \bigl(\Flow_{\dt/M}\bigr)^{\circ M},\qquad M = 32 .
\]
\end{definition}

\begin{remark}[The central methodological confound]
\label{rem:tautology}
This is not a numerical nicety; it is what stops the study from being circular. Model C
is, by construction, \emph{a single Strang step with learned rates}. If the training
targets were produced by a single Strang step at the same $\dt$, the true one-step map
would lie \emph{inside} Model C's hypothesis class, and ``structure beats FNO'' would be
a tautology --- C would win because it was handed the answer, not because structure
helps. Targets must therefore come from a converged, substepped reference, and every
error is quoted relative to it.
\end{remark}

\begin{definition}[The splitting floor]
\label{def:epssplit}
\[
\epssplit \;=\; \bigl\| \Flow_{\dt} - \Flow^{\mathrm{ref}}_{\dt}\bigr\|_{\text{rel }L^2},
\]
the relative $L^2$ error of one Strang step at $\dt$ against the substepped reference,
measured on the production distribution.
\end{definition}

$\epssplit$ is an \emph{analytic lower bound on the one-step error of any model whose
hypothesis class is a single split step at $\dt$ with the exact rates}. Its role in the
argument is precise:

\begin{itemize}[leftmargin=1.3em,itemsep=0.2em]
\item A structured model that trains \emph{down to} $\epssplit$ has recovered the physics
--- there is nothing left in its class to find.
\item A black-box model that goes \emph{below} $\epssplit$ is exploiting non-split
structure that the splitting cannot represent. That is a finding about expressivity, not
a failure of the bound.
\item A black-box model that sits well \emph{above} $\epssplit$ has left headroom for the
structured models. This is what was measured: the FNO is $11\times$ above the floor.
\end{itemize}

\subsection{The fitted floor: the bound a \emph{learned} split step actually faces}
\label{sec:fittedfloor}

Definition~\ref{def:epssplit} measures the error of the Strang map built from the
\emph{true} rates $\kappa = -\alpha|k|^2$, $\nu = \beta\rho - V$. A learned split step is
under no obligation to use those. It may pick \emph{effective} rates that absorb part of
the $O(\dt^2)$ commutator, so the quantity that genuinely lower-bounds Model C is
\begin{equation}
\label{eq:fittedfloor}
\epssplit^{\mathrm{fit}}
\;=\;
\inf_{\theta}\;\bigl\|\Flow^\theta_{\dt} - \Flow^{\mathrm{ref}}_{\dt}\bigr\| .
\end{equation}

Two nested families are fitted, both shared across the whole batch (a single model must
serve every $(\alpha,\beta)$):
\[
\text{(scalar)}\quad \kappa = -c_0\,\alpha|k|^2,\quad \nu = c_1(\beta\rho - V);
\qquad
\text{(per-mode)}\quad \kappa = -\alpha|k|^2\,g(|k|^2),
\]
with $g$ a free vector initialised at $1$. The per-mode family is exactly the \code{K2}
parameterisation Phase~4 gives Model C1, so its achieved error is the target C1-K2 must
be measured against. Note that $g$ \emph{multiplies} $\alpha|k|^2$ rather than replacing
it: a free function of $|k|^2$ alone could not represent the required $\alpha$-scaling.

\begin{observation}[$\epssplit$ is a genuine bound]
\label{obs:floor}
Fitting the effective generators recovers the exact rates to within $10^{-4}$
($c_0 = 1.0000008$, $c_1 = 0.9999979$) and improves the error by at most $\sim$5\%:
$\epssplit^{\mathrm{fit}}/\epssplit = 0.9999978$ for the scalar family and $0.99999$ for
the free 32-parameter per-mode family. \emph{The Strang commutator is very nearly
orthogonal to the span of the generators.} There is no ``cheap'' effective rate that
absorbs it.
\end{observation}

\begin{observation}[Reference self-convergence]
$M{=}32$ versus $M{=}64$ differ by $1.50\times10^{-7}$, i.e.\ $0.24\%$ of $\epssplit$.
The reference is converged well below the quantity it is used to measure.
\end{observation}

\sourcenote{src/spno/solvers/split\_step.py}
\newpage
\part*{Part II \quad What each model class guarantees}
\addcontentsline{toc}{part}{Part II \; What each model class guarantees}

\section{Model A --- the unrestricted Fourier neural operator}
\label{sec:fno}

\subsection{The spectral convolution}

\begin{definition}[Mode-truncated global convolution]
For a hidden state $u\in\R^{C\times N}$ and learned complex weights
$R\in\C^{C\times C'\times m}$,
\[
\bigl(\mathcal{S}u\bigr) \;=\; \mathcal{F}^{-1}\Bigl( \mathbf{1}_{|k|<m}\cdot R_k\,\widehat{u}_k \Bigr),
\]
where $\mathcal{F}$ is the real FFT and $m$ is the retained mode budget
(\code{modes}$=16$). Modes above $m$ are set to zero, not merely damped.
\end{definition}

An FNO layer is $u\mapsto \sigma\bigl(\mathcal{S}u + Wu\bigr)$ with $W$ a pointwise
($1\times1$) convolution. The full model lifts $5$ real input channels
$(\Real\psi,\Imag\psi,V,\alpha,\beta)$ to width $64$, applies four such layers, and
projects to $2$ output channels read as $(\Real\psi',\Imag\psi')$.

\begin{proposition}[Grid independence, and its exact scope]
\label{prop:fnogrid}
The spectral weights $R$ are indexed by \emph{frequency}, and the inverse transform
receives the target length $n$ at call time. Hence the same weights define the same
operator on any grid --- \textbf{for $|k|\le m$ only}. They say nothing whatever about
modes above the budget.
\end{proposition}

This is the \emph{only} source of the FNO's much-advertised resolution transfer, and
\S\ref{sec:resolution} is largely about not overclaiming it.

\begin{guarantee}[Model A guarantees nothing]
\label{guar:fnonone}
$\Flow^{\mathrm{A}}$ is an arbitrary trained map $\C^N\to\C^N$. It does not conserve
$\Mass_h$, does not conserve $\Ham$, is not time-reversible, is not $U(1)$-equivariant,
and does not respect \eqref{eq:dispersion}. Any agreement with these is learned and
approximate.
\end{guarantee}

\begin{remark}[$\dt$ does not enter the map]
\label{rem:fnodt}
The FNO ignores its $\dt$ argument entirely. Consequently:
(i) evaluating it at $-\dt$ re-applies the \emph{forward} map, so any ``reversibility''
number measured that way is an artefact --- \code{evaluate\_reversibility} refuses rather
than returning it; and (ii) it has no basis for $\dt$-transfer. The base class
encodes this as \code{supports\_time\_reversal} and \code{supports\_dt\_transfer} flags
that raise rather than silently returning a wrong answer.
\end{remark}

\subsection{Three design choices, each tied to a confound}

\begin{enumerate}[leftmargin=1.4em,itemsep=0.35em]
\item \textbf{Parameter conditioning is rescaled to $O(1)$.} $\alpha$ and $\beta$ are
broadcast over the grid and mapped affinely from their sampling ranges to $[-1,1]$:
\[
\tilde\alpha = \frac{2(\alpha-\alpha_{\min})}{\alpha_{\max}-\alpha_{\min}} - 1 .
\]
A consequence worth recording: with $\alpha\in[0.7,1.1]$, the value $\alpha=0$ enters the
lift channel at $\tilde\alpha = -4.5$. Any $\alpha$-continuation experiment
(\S\ref{sec:continuation}) that anchors near $\alpha=0$ is therefore probing
$\alpha$-extrapolation, and a failure there says nothing about whether the
$\alpha$-derivative information was extracted.

\item \textbf{No absolute-coordinate channel.} The joint operator $(\psi,V)\mapsto\psi'$
is translation-equivariant --- the potential, not the coordinate, is what breaks
homogeneity. Feeding $x$ into the lift lets the model memorise position instead of
learning the operator. Kept behind a flag as an ablation.

\item \textbf{A single global amplitude scale, not per-sample normalisation.} The
nonlinearity $\beta|\psi|^2\psi$ is \emph{not} scale-invariant, and mass varies across
samples by design (\S\ref{sec:datadesign}). Normalising each sample by its own norm would
destroy exactly the information the nonlinear term depends on.
\end{enumerate}

\subsection{The three separable causes of high-$k$ failure}
\label{sec:highk}

When an FNO does badly at large $|k|$, three mechanisms are in play, and only one of them
is interesting. Keeping them apart is the entire design of experiment G7.

\begin{enumerate}[label=(\alph*),leftmargin=1.6em,itemsep=0.25em]
\item \textbf{Hard mode truncation beyond $m$.} The spectral branch contributes nothing
above $|k|=m=16$. High-$k$ content can only reach the output through the pointwise
branches and the nonlinearity. \emph{Trivial}, and removable by matching $m$ to Nyquist.
\item \textbf{The oscillatory dependence of $\eu^{-\iu\alpha|k|^2\dt}$ on $\alpha$.} The
FNO must route this through a pointwise lift of a broadcast scalar. \emph{This is the
interesting one} --- it is a statement about inductive bias.
\item \textbf{No training energy above $\ktrain$.} Initial conditions are band-limited to
$|k|\le 8$; nothing above that is supervised. \emph{No architecture fixes this.}
\end{enumerate}

\sourcenote{src/spno/models/fno.py}

\section{Model B --- hard constraint projection}
\label{sec:projection}

\begin{definition}[Mass projection]
\label{def:proj}
\[
\Pi_{M}(\psi) \;=\; \sqrt{\frac{M}{\Mass_h[\psi]}}\;\psi,
\qquad
\Flow^{\mathrm{B}}_{\dt}(\psi_n) \;=\; \Pi_{\Mass_h[\psi_n]}\bigl(\Flow^{\mathrm{A}}_{\dt}(\psi_n)\bigr).
\]
\end{definition}

\begin{proposition}[$\Pi_M$ is the metric projection onto the mass sphere]
\label{prop:metricproj}
Let $S_R=\{u:\|u\|_{L^2_h}=R\}$ with $R=\sqrt{M}$. For $\psi\neq 0$,
\[
\Pi_M(\psi) \;=\; \argmin_{u\in S_R}\|u-\psi\|_{L^2_h},
\qquad\text{and}\qquad
\min_{u\in S_R}\|u-\psi\| \;=\; \bigl|\,R - \|\psi\|\,\bigr| .
\]
\end{proposition}

\begin{proof}
For $u\in S_R$, $\|u-\psi\|^2 = R^2 + \|\psi\|^2 - 2\Real\langle u,\psi\rangle$, so
minimising the distance maximises $\Real\langle u,\psi\rangle$. By Cauchy--Schwarz,
$\Real\langle u,\psi\rangle \le \|u\|\|\psi\| = R\|\psi\|$ with equality iff $u$ is a
positive multiple of $\psi$; the unique such $u$ in $S_R$ is $R\psi/\|\psi\|$, which is
$\Pi_M(\psi)$. Substituting gives $\|u-\psi\|^2 = (R-\|\psi\|)^2$.
\end{proof}

So the name is earned: this is genuinely the closest mass-correct field, not merely a
convenient rescaling.

\begin{guarantee}[Exact mass, and nothing else]
\label{guar:B}
$\Mass_h[\Flow^{\mathrm{B}}_{\dt}\psi_n] = \Mass_h[\psi_n]$ exactly in exact arithmetic,
for every parameter value. Model B guarantees \textbf{no other} property: not phase, not
energy, not reversibility, not the dispersion relation.
\end{guarantee}

\begin{remark}[Two training modes, and they are not equivalent]
\emph{Projection in the loop} makes $\Pi$ part of the forward pass, so gradients flow
through it and the core learns to produce fields whose \emph{shape} is right with
amplitude handled downstream. \emph{Post hoc} trains the core unconstrained and applies
$\Pi$ only at evaluation, so the core still spends capacity on the amplitude. They must
be reported separately.
\end{remark}

\begin{observation}[What one hard invariant actually buys]
\label{obs:B}
Across $3$ seeds at matched budget: the projection removes mass drift across
\textbf{13 orders of magnitude} ($1.61\times10^{-2}\to1.3\times10^{-15}$ at step 100) and
buys \textbf{$\sim$12\% rollout accuracy} ($2.77\times10^{-2}\to2.44\times10^{-2}$),
of which $\sim$7 points come from the projection alone and $\sim$5 from training with it.
Energy remains \textbf{secular} for A, B-post and B-loop alike, in $9/9$ seed--model runs.
\emph{Projecting mass does not fix energy.} If that reads as a disappointing result, it
is the intended lesson: one scalar constraint is one scalar constraint.
\end{observation}

\sourcenote{src/spno/models/projected.py, src/spno/domain.py}

\section{Model C --- structure encoded in the architecture}
\label{sec:learnedsplit}

\subsection{The parameterisation}

\begin{definition}[Learned split step]
\label{def:learnedsplit}
\[
\boxed{\;
\Flow^{\theta}_{\dt} \;=\; \Kop^{\theta}_{\dt/2}\circ\Nop^{\theta}_{\dt}\circ\Kop^{\theta}_{\dt/2}
\;}
\]
with
\[
\widehat{\Kop^{\theta}_{\tau}\psi}(k) \;=\; \eu^{\,\iu\tau\,\kappa_\theta(|k|^2,\alpha,\beta)}\;\widehat\psi(k),
\qquad
\Nop^{\theta}_{\tau}\psi \;=\; \eu^{\,\iu\tau\,\nu_\theta(\rho,V,\alpha,\beta)}\;\psi,
\qquad \rho=|\psi|^2,
\]
and \textbf{both $\kappa_\theta$ and $\nu_\theta$ real-valued}. The true values are
$\kappa = -\alpha|k|^2$ and $\nu = \beta\rho - V$.
\end{definition}

Everything that follows depends on exactly two facts: the rates are \emph{real}, and
$\nu_\theta$ sees $\psi$ \emph{only through $\rho$}.

\subsection{The main theorem}

\begin{theorem}[Structure of the learned split step]
\label{thm:learnedmain}
For every $\theta$ --- untrained weights included --- the map $\Flow^\theta_{\dt}$ of
Definition~\ref{def:learnedsplit} satisfies:
\begin{enumerate}[label=\textup{(\arabic*)},leftmargin=2.2em,itemsep=0.3em]
\item \textbf{Exact mass.} $\Mass_h[\Flow^\theta\psi] = \Mass_h[\psi]$.
\item \textbf{Exact reversibility.} $\Flow^\theta_{-\dt}\circ\Flow^\theta_{\dt} = \Id$.
\item \textbf{Symplecticity, as the Strang splitting of a learned Hamiltonian.} There
exists $\Ham_\theta = \Ham_{K,\theta} + \Ham_{N,\theta}$ with
\[
\Ham_{K,\theta} = -\!\int \kappa_\theta(|k|^2)\,|\widehat\psi(k)|^2,
\qquad
\Ham_{N,\theta} = -\!\int F_\theta(\rho,V),\quad \frac{\partial F_\theta}{\partial\rho}=\nu_\theta,
\]
such that each substep is the \emph{exact} flow of its sub-Hamiltonian. Hence
$\Flow^\theta_{\dt}$ is symplectic and time-symmetric.
\item \textbf{Global $U(1)$ equivariance.}
$\Flow^\theta(\eu^{\iu c}\psi) = \eu^{\iu c}\,\Flow^\theta(\psi)$ for every $c\in\R$.
\end{enumerate}
\end{theorem}

\begin{proof}\leavevmode

\emph{(1).} $\nu_\theta$ is real-valued by construction, so $\eu^{\iu\dt\nu_\theta}$ has
modulus one pointwise and $|\psi_j|$ is unchanged at every grid point. $\kappa_\theta$ is
real-valued, so $\eu^{\iu\tau\kappa_\theta}$ has modulus one mode by mode and
Lemma~\ref{lem:parseval} applies. Composition of norm-preserving maps.

\emph{(2).} The kinetic multipliers cancel in pairs exactly as in
Theorem~\ref{thm:strangrev}, because $\kappa_\theta$ does not depend on $\psi$ at all ---
it is a function of $(|k|^2,\alpha,\beta)$ only, so the forward and backward half-steps
use \emph{identical} rates. For the local step, let $w$ be its input, $\rho_w=|w|^2$, and
$v = \eu^{\iu\dt\nu_\theta(\rho_w,\cdot)}w$. By (1) applied pointwise, $|v|^2 = \rho_w$,
so $\nu_\theta$ evaluated at $v$ equals $\nu_\theta$ evaluated at $w$ --- the backward
step reads the same phase field --- and
$\eu^{-\iu\dt\nu_\theta(\rho_w,\cdot)}v = w$. This is the step where $\rho$-only
dependence is used, and it is the step that fails for C3.

\emph{(3).} Consider first the kinetic generator. The flow
$\partial_t\widehat\psi_k = \iu\kappa_k\widehat\psi_k$ is Hamiltonian with
$\Ham_{K,\theta} = -\sum_k \kappa_k|\widehat\psi_k|^2$ (up to the DFT normalisation),
since
\[
\frac{\partial \Ham_{K,\theta}}{\partial\overline{\widehat\psi_k}} = -\kappa_k\widehat\psi_k
\quad\Longrightarrow\quad
\iu\,\partial_t\widehat\psi_k = \frac{\partial\Ham_{K,\theta}}{\partial\overline{\widehat\psi_k}}
\quad\Longleftrightarrow\quad
\partial_t\widehat\psi_k = \iu\kappa_k\widehat\psi_k .
\]
$\Ham_{K,\theta}$ is real because $\kappa_\theta$ is.

Now the local generator $\partial_t\psi = \iu\nu_\theta(\rho,V)\psi$. Fix $(V,\alpha,\beta)$
and define
\[
F_\theta(\rho) \;=\; \int_0^{\rho}\nu_\theta(s,V,\alpha,\beta)\dd s,
\]
which exists because $\nu_\theta$ is a continuous (indeed smooth) function of the
\emph{scalar} $\rho$ --- this is the point at which the pointwise structure is
indispensable, and it is why an antiderivative always exists no matter what the network
has learned. Set $\Ham_{N,\theta} = -\int F_\theta(\rho)\dd x$. Using
$\rho = \psi\bar\psi$ so that $\partial\rho/\partial\bar\psi = \psi$,
\[
\frac{\delta \Ham_{N,\theta}}{\delta\bar\psi}
= -\frac{\partial F_\theta}{\partial\rho}\cdot\frac{\partial\rho}{\partial\bar\psi}
= -\nu_\theta\,\psi
\quad\Longrightarrow\quad
\iu\,\partial_t\psi = -\nu_\theta\psi \cdot(-1)
\quad\Longleftrightarrow\quad
\partial_t\psi = \iu\nu_\theta\psi ,
\]
as required. Finally, by the argument of Lemma~\ref{lem:local} (which used only that
$\nu_\theta$ is real), $\rho$ is constant along this flow, so the closed-form multiplier
$\eu^{\iu\tau\nu_\theta}$ \emph{is} the exact time-$\tau$ flow --- not an approximation to
it. Each substep is therefore an exact Hamiltonian flow, hence symplectic;
$\Flow^\theta_{\dt}$ is a composition of symplectic maps, and is time-symmetric by the
palindromic ordering together with (2).

\emph{(4).} $\Kop^\theta$ is linear, so it commutes with multiplication by the constant
$\eu^{\iu c}$. And $\rho$ is phase-blind: $|\eu^{\iu c}\psi|^2 = |\psi|^2$. Hence
$\nu_\theta$ is unchanged and $\Nop^\theta$, being multiplication by a $\psi$-independent
(given $\rho$) pointwise factor, also commutes with the global phase. Composing,
$\Flow^\theta(\eu^{\iu c}\psi) = \eu^{\iu c}\Flow^\theta(\psi)$.
\end{proof}

\begin{remark}[Two routes to (1)]
Part (1) also follows from (3) and (4) by Noether's theorem: $\Ham_\theta$ depends on
$\psi$ only through $|\widehat\psi|^2$ and $\rho$, so it is $U(1)$-invariant, and the
conserved charge of that symmetry is $\Mass_h$. The direct modulus-one argument is the
one the tests assert, because it holds mode by mode and grid point by grid point and is
therefore checkable to machine precision.
\end{remark}

\subsection{The sharp, falsifiable prediction}
\label{sec:sharpprediction}

Part (3) is not decoration. It is the claim the study can be wrong about.

\begin{corollary}[Bounded, not secular, energy error]
\label{cor:boundedenergy}
$\Flow^\theta$ does \emph{not} conserve the true Hamiltonian $\Ham$. What it does is
integrate a \emph{learned} Hamiltonian $\Ham_\theta$ exactly-symplectically. By backward
error analysis (Corollary~\ref{cor:modham} applied to $\Ham_\theta$), the drift in
$\Ham_\theta$ is bounded; the drift in $\Ham$ is therefore governed by the fixed
model-error term $\|\Ham_\theta - \Ham\|$ evaluated along the trajectory. That is a
constant of the model, not a function of the horizon. \textbf{Prediction: energy drift is
bounded and oscillatory, not secular.}
\end{corollary}

Phases 2--3 measured \emph{secular} energy drift for the FNO and for the mass-projected
FNO in $9/9$ seed--model runs (Observation~\ref{obs:B}), so this prediction is falsifiable
against a real baseline rather than against a straw man.

\begin{remark}[The reversible-KAM caveat --- state it before the result comes in]
\label{rem:kam}
Symmetric \emph{and} reversible methods already enjoy bounded-energy results for
reversible systems, independently of symplecticity. Model C2 is reversible but not
symplectic (\S\ref{sec:c2}). It is therefore \textbf{unlikely that C1 and C2 separate on
energy drift}, and a finding that they do not is not evidence against
Corollary~\ref{cor:boundedenergy}. Frame the C1-vs-C2 comparison as \emph{sample
efficiency versus expressivity}, not as a test of symplecticity.
\end{remark}

\subsection{C2 --- what survives a nonlocal phase}
\label{sec:c2}

Model C2 replaces the pointwise $\nu_\theta$ with an FNO reading the four real channels
$(\rho, V, \alpha, \beta)$. The phase is now a \emph{nonlocal} functional of $\rho$.

\begin{guarantee}[C2]
\label{guar:c2}
Parts (1), (2) and (4) of Theorem~\ref{thm:learnedmain} hold verbatim. Part (3) does not.
\end{guarantee}

\begin{proof}
(1) and (4) used only that $\nu_\theta$ is real and phase-blind, both of which survive.
(2) used that $\rho$ is unchanged by the local step --- and it is unchanged
\emph{as an entire field}, not merely at the point being updated. A nonlocal $\nu$ reads
the whole $\rho$ field, which is frozen, so the backward step still sees the identical
phase field and still cancels exactly.

For (3), the local generator $\partial_t\psi = \iu\nu[\rho]\psi$ is Hamiltonian iff
$\nu$ is the functional derivative of some $G[\rho]$, i.e.\ iff its linearisation is
self-adjoint:
\[
\frac{\delta\nu(x)}{\delta\rho(y)} \;=\; \frac{\delta\nu(y)}{\delta\rho(x)}
\qquad\text{for all }x,y .
\]
This is a codimension-large symmetry condition on the operator, and an unconstrained FNO
does not satisfy it. Hence C2's local substep is generically not the flow of a
sub-Hamiltonian and the composition is not symplectic.
\end{proof}

\begin{remark}[C2 is the fair fight]
Model A and Model C2 share the \emph{same backbone}, the same hyperparameters and
essentially the same parameter count ($287{,}746$ vs $288{,}770$, a ratio of
$1.004\times$). They differ in exactly one thing: whether the network's output \emph{is}
the field, or is a \emph{phase applied to} the field. That is the cleanest available
isolation of ``structure'' from ``capacity''. C1, at $2{,}466$ parameters, is $117\times$
smaller than A --- which is why parameter counts must appear in every table.
\end{remark}

\subsection{C3 --- the control, and why ``irreversible'' would be the wrong word}
\label{sec:c3}

Model C3 feeds $\Real\psi$ and $\Imag\psi$ to the phase network instead of $\rho$.

\begin{guarantee}[C3]
\label{guar:c3}
Part (1) of Theorem~\ref{thm:learnedmain} still holds --- the multiplier is still modulus
one. Part (4) is destroyed outright. Part (2) degrades from \emph{exact} to
\emph{second order}.
\end{guarantee}

\begin{proof}[Proof of the reversibility statement]
Write the local step as $u\mapsto v = \eu^{\iu\dt\,\nu(u)}u$, where now $\nu$ depends on
$u$ and not merely on $|u|$. The backward step gives
\[
\Nop_{-\dt}(v) \;=\; \eu^{-\iu\dt\,\nu(v)}v \;=\; \eu^{\,\iu\dt\,[\nu(u)-\nu(v)]}\,u .
\]
Now $v = u + O(\dt)$, indeed $v = u + \iu\dt\nu(u)u + O(\dt^2)$, so by smoothness of
$\nu$,
\[
\nu(v) - \nu(u) \;=\; D\nu[u]\bigl(\iu\dt\,\nu(u)u\bigr) + O(\dt^2) \;=\; O(\dt).
\]
Hence the residual phase is $\dt\cdot O(\dt) = O(\dt^2)$, and it is generically nonzero,
because $D\nu[u]$ applied to $\iu\nu(u)u$ --- a purely ``rotational'' perturbation ---
vanishes only if $\nu$ happens to be phase-blind, which is exactly what C3 gives up.
Since the kinetic halves are unitary and cancel exactly, the round-trip error of the full
step is $O(\dt^2)$.
\end{proof}

\begin{observation}[Measured order $2.00$]
\label{obs:c3order}
At untrained weights, C3's forward--backward violation is $\approx1\times10^{-9}$ at
$\dt=0.01$, against $\approx7\times10^{-16}$ and \emph{$\dt$-independent} for C1/C2. Its
measured $\dt$-scaling is order $2.00$.
\end{observation}

\begin{trap}[Absolute thresholds hide power laws]
\label{trap:thresholds}
$10^{-9}$ is small enough that any fixed tolerance would call C3 ``reversible''. It is
not; it is second-order. Regimes are therefore classified by \emph{measured $\dt$-scaling}
--- $\log_2(e_{i+1}/e_i)$ compared against an expected order --- never by an absolute
threshold. The defensible claim is:
\begin{center}
\emph{$\rho$-only parameter sharing makes reversibility \textbf{exact rather than
second-order}}
\end{center}
and \emph{not} ``C3 is irreversible''.
\end{trap}

C3 has two jobs. It is the direct control for Model B --- both preserve mass and nothing
else, at very different capacities --- and it is the \textbf{paired negative} proving the
C1/C2 reversibility and $U(1)$ tests are not vacuous. A conservation test that would pass
on a model with no such structure tests nothing.

\subsection{The rate ladders: approximation versus optimisation}
\label{sec:ladders}

Both learned rates are, in truth, \emph{products}: $\kappa = -\alpha\cdot|k|^2$ and
$\nu = \beta\cdot\rho - V$. A $\tanh$ MLP handed $(|k|^2,\alpha,\beta)$ or $(\rho,V)$ as
features must \emph{discover} the bilinear form. Whether it can is a question about
learning, not about expressivity, and the two must not be conflated. Hence two ladders:

\begin{center}
\begin{tabular}{@{}llp{0.52\linewidth}@{}}
\toprule
\textbf{rung} & \textbf{kinetic $\kappa$} & \textbf{what is handed to the network}\\
\midrule
\code{K0} & $\mathrm{MLP}(\widetilde{k^2},\alpha,\beta)$ & nothing --- must find the product\\
\code{K1} & $\mathrm{MLP}(\widetilde{k^2},\,\alpha\widetilde{k^2},\,\alpha,\beta)$ & the product as an extra feature\\
\code{K2} & $-\alpha|k|^2\bigl(1+\mathrm{MLP}(\cdot)\bigr)$ & the functional form; MLP only corrects\\
\midrule
\textbf{rung} & \textbf{local $\nu$} & \\
\midrule
\code{L0} & $\mathrm{MLP}(\rho,V;\alpha,\beta)$ & nothing\\
\code{L1} & $\mathrm{MLP}(\rho,V,\beta\rho;\alpha,\beta)$ & the product $\beta\rho$\\
\code{L2} & $\beta\rho - V + \mathrm{MLP}(\cdot)$ & the exact rate; MLP only corrects\\
\bottomrule
\end{tabular}
\end{center}

\code{K2} and \code{L2} are zero-initialised at the last layer, so they \emph{start} at
the exact rate.

\begin{remark}[Why \code{K0}/\code{L0} are the headline]
A win above $\kwrap$ at \code{K2} is uninterpretable: the model was told the dispersion
relation. A win at \code{K0} means it \emph{learned} it. The ladder is reported in full,
because if \code{L0} cannot reach the floor that is itself a finding --- about how hard
the learning problem is, distinct from how expressive the class is.
\end{remark}

\begin{observation}[The approximation floor of a free local rate]
\label{obs:ladder}
Measured on the production distribution, a free MLP plateaus at $2.1\%$ relative RMSE on
$\nu$ (a wider, deeper one at $2.6\%$ --- an \emph{optimisation} limit, not a capacity
one). End-to-end, with the kinetic rate held exact:
\[
\code{L0}: 1.18\,\epssplit,\qquad
\code{L1}: 1.03\,\epssplit,\qquad
\code{L2}: 0.93\,\epssplit .
\]
So a free local rate inflates C1's one-step error by $\sim$18\%, \emph{not} by the factor
of $\sim$3 that a naive $\dt\cdot|\nu|\cdot e$ bound predicts --- the phase error is partly
incoherent across the grid and does not accumulate as that bound assumes. \code{L2}
landing slightly \emph{below} the floor is expected, and matches
Observation~\ref{obs:floor}: a free correction to the generators absorbs a few percent of
the Strang commutator.
\end{observation}

\begin{trap}[Normalise $|k|^2$ before it reaches a $\tanh$]
At $N=64$, $|k|^2$ reaches $1024$ and saturates every $\tanh$ in the network. The
kinetic MLP therefore consumes $\widetilde{k^2} = |k|^2/\max|k|^2$. The normaliser is
\textbf{frozen at the training grid} when the model is rebound to a finer grid --- see
Proposition~\ref{prop:frozennorm}.
\end{trap}

\subsection{Summary of guarantees}

\begin{center}
\begin{tabular}{@{}lcccc@{}}
\toprule
& \textbf{mass} & \textbf{reversibility} & \textbf{symplectic} & \textbf{$U(1)$}\\
\midrule
A \; (FNO) & --- & --- & --- & ---\\
B \; (FNO $+$ projection) & exact & --- & --- & ---\\
C1 (pointwise $\rho$ phase) & exact & exact & exact & exact\\
C2 (FNO-on-$\rho$ phase) & exact & exact & --- & exact\\
C3 (Re/Im $\psi$ phase) & exact & order $2$ & --- & ---\\
\midrule
reference (Strang) & exact & exact & exact & exact\\
\bottomrule
\end{tabular}
\end{center}

Every ``exact'' in this table is asserted in the test suite \textbf{at random untrained
weights}, each with a paired negative demonstrating the test can fail. Positives bound the
worst case ($\max < 10^{-13}$); negatives bound the best case ($\min > 10^{-6}$).

\sourcenote{src/spno/models/split\_learned.py}
\newpage
\part*{Part III \quad Spectral identifiability --- the thesis spine}
\addcontentsline{toc}{part}{Part III \; Spectral identifiability}

\section{The one-step multiplier and what it can tell you}
\label{sec:identifiability}

The headline question of the study is \emph{architecture-relative identifiability}:

\begin{quote}
Single-$\alpha$ supervision makes $\omega(k)$ unrecoverable above
$\kwrap = \sqrt{\pi/(\alpha\dt)}$ for \textbf{every} model --- a clean theorem.
$\alpha$-varying supervision makes it recoverable \textbf{in principle}, because
$\partial_\alpha\arg m(k,\alpha) = -k^2\dt$ is wrap-free. The question is which hypothesis
classes actually extract it. \emph{The data contains it; only the right inductive bias
gets it out.}
\end{quote}

This section proves both halves.

\subsection{The probe}

\begin{definition}[One-step multiplier]
\label{def:multiplier}
Fix a constant potential $V\equiv V_0$, an amplitude $A$, and parameters $(\alpha,\beta)$.
For any one-step map $\Flow$ (a trained model or the reference solver), let
\[
m(k) \;=\; \frac{\widehat{\Flow(\psi)}(k)}{\widehat{\psi}(k)},
\qquad \psi = A\,\eu^{\iu k\cdot x}.
\]
\end{definition}

\begin{proposition}[The splitting is \emph{exact} on plane waves in a constant potential]
\label{prop:strangexactpw}
For $V\equiv V_0$ and $\psi = A\eu^{\iu kx}$, the single Strang step, the substepped
reference, and the exact flow all coincide:
\[
\Flow_{\dt}\psi \;=\; \Flow^{\mathrm{ref}}_{\dt}\psi \;=\; \eu^{-\iu\omega\dt}\,\psi,
\qquad \omega = \alpha|k|^2 - \beta A^2 + V_0 .
\]
\end{proposition}

\begin{proof}
$\Kop_{\dt/2}\psi = \eu^{-\iu\alpha|k|^2\dt/2}\psi$, a \emph{constant} multiple of $\psi$,
so the midpoint state has $\rho \equiv A^2$, constant in $x$. The local phase
$\eu^{\iu(\beta A^2-V_0)\dt}$ is therefore also a constant multiple. Composing the three
constants gives $\eu^{-\iu(\alpha|k|^2 - \beta A^2 + V_0)\dt}$. The same argument applies
to each substep of $\Flow^{\mathrm{ref}}$, and the constants compose to the same total.
By Proposition~\ref{prop:planewave} this is the exact flow.
\end{proof}

\begin{remark}[The floor for probe residuals is roundoff, not $\epssplit$]
\label{rem:proberesidualfloor}
Proposition~\ref{prop:strangexactpw} has a sharp consequence that is easy to get wrong:
$\epssplit = 6.239\times10^{-5}$ was measured on \emph{random band-limited fields}, where
$\rho$ varies in $x$ and the commutator is nonzero. It is \textbf{not} the floor for
dispersion-probe residuals. On plane waves in a constant potential the floor is
\code{float64} roundoff, $\sim10^{-13}$ in $\omega$ at $\dt=0.01$, and the probe
self-validation gate is set accordingly at $10^{-10}$.
\end{remark}

\begin{remark}[What the probe is \emph{not}]
Three limits, each load-bearing and each stated in the thesis rather than buried:
\begin{enumerate}[label=(\roman*),leftmargin=1.8em,itemsep=0.2em]
\item Plane waves are \textbf{out of distribution} for every model here --- all are
trained on random band-limited fields. This probes the \emph{learned operator}, not
generalisation.
\item The probe is valid for \textbf{constant $V_0$ only}, because that is the only case
in which a plane wave is an exact solution.
\item The branch lift (\S\ref{sec:wraptheorem}) selects the branch \emph{using the truth}.
Below $\kwrap$ the branch index is zero and the lift is a consistency check; above it the
lift is an assist and the ``recovered'' $\omega$ is not an independent measurement. That
asymmetry \emph{is} the G5a result, not a defect in the instrument.
\end{enumerate}
\end{remark}

\subsection{The probe amplitude must not smuggle in a mass shift}
\label{sec:probeamp}

By \eqref{eq:dispersion}, $\omega$ depends on the amplitude through $\beta A^2$; by
Proposition~\ref{prop:pwmass}, the plane wave's mass is $A^2\prod_j L_j$. A probe
amplitude outside the training mass distribution therefore makes every high-$k$ reading a
measurement of \emph{two} shifts at once --- spectral and mass --- with no way to
attribute the failure.

\begin{definition}[Probe amplitude]
With the training mass range $[M_-,M_+]$ and $\overline{M}=\tfrac12(M_-+M_+)$,
\[
A \;=\; \sqrt{\overline{M}\,\Big/\!\prod_j L_j} .
\]
For the production configuration, $M\in[1,3]$ and $\prod L_j = 2\pi$, so
$A = \sqrt{2/(2\pi)} \approx 0.5642$.
\end{definition}

\subsection{The negative theorem}
\label{sec:wraptheorem}

\begin{theorem}[Single-$\alpha$ ambiguity]
\label{thm:wrap}
Let the supervision consist of one-step data at a \textbf{fixed} $\alpha$ and fixed
$\dt$. Then for each mode $k$, the data determine $\omega(k)$ only modulo $2\pi/\dt$.
Writing the principal branch
\[
\omega_{\mathrm{princ}}(k) \;=\; -\frac{\arg m(k)}{\dt} \;\in\; \Bigl[-\frac{\pi}{\dt},\,\frac{\pi}{\dt}\Bigr),
\]
one has $\omega_{\mathrm{princ}}(k) = \omega(k)$ if and only if $|\omega(k)|\dt \le\pi$.
In particular, with the kinetic term dominant, the horizon is
\begin{equation}
\label{eq:kwrap}
\boxed{\;\kwrap \;=\; \sqrt{\frac{\pi}{\alpha\,\dt}}\;}
\end{equation}
and above it the branch is \textbf{genuinely unrecoverable} from the data --- for every
model, of every architecture.
\end{theorem}

\begin{proof}
The only thing a one-step map on a plane wave reveals about mode $k$ is the complex
number $m(k)$, and by Proposition~\ref{prop:strangexactpw} the ground truth is
$m(k)=\eu^{-\iu\omega(k)\dt}$. The map $\omega\mapsto\eu^{-\iu\omega\dt}$ is
$2\pi/\dt$-periodic, so $m(k)$ determines $\omega(k)$ exactly up to that period. The
principal-branch statement is the definition of $\arg$ taking values in $(-\pi,\pi]$.
Setting $|\omega|\dt=\pi$ with $\omega\approx\alpha k^2$ and solving for $k$ gives
\eqref{eq:kwrap}.
\end{proof}

\begin{corollary}[Multi-step data at the same $\dt$ adds nothing]
\label{cor:multistep}
Let $m_n(k)$ be the $n$-step multiplier. Then $m_n = m^n$, so $m_n$ is a deterministic
function of $m$: the $n$-step observation is measurable with respect to the $\sigma$-algebra
generated by the one-step observation and carries no additional information about
$\omega$. (Conversely $m_n$ determines $m$ only up to an $n$-th root, so it is strictly
weaker.)
\end{corollary}

\begin{proof}
A plane wave in a constant potential remains a plane wave of the same amplitude under the
flow (Proposition~\ref{prop:strangexactpw}), so iterating multiplies the coefficient by
$m$ each time.
\end{proof}

\subsection{The alias family, constructed}
\label{sec:alias}

Theorem~\ref{thm:wrap} has a constructive shadow that makes it testable.

\begin{proposition}[Indistinguishable $\alpha$ values at a given mode]
\label{prop:alias}
Fix $k$ and $\dt$. Define the alias family
\begin{equation}
\label{eq:aliasfamily}
\alpha_n \;=\; \alpha_0 \;+\; \frac{2\pi n}{k^2\dt},\qquad n\in\Z .
\end{equation}
Then $m(k;\alpha_n) = m(k;\alpha_0)$ for every $n$: the one-step map restricted to mode
$k$ is \textbf{identical} across the family.
\end{proposition}

\begin{proof}
$\arg m(k;\alpha) = -(\alpha k^2 - \beta A^2 + V_0)\dt$. Replacing $\alpha_0$ by
$\alpha_0 + 2\pi n/(k^2\dt)$ changes the argument by $-2\pi n$, i.e.\ not at all as a
complex number.
\end{proof}

\begin{observation}[The computed family at $k=30$]
\label{obs:alias}
With $\dt=0.01$, $k=30$: $2\pi/(k^2\dt) = 2\pi/9 = 0.69813$, so from $\alpha_0=0.9$ the
family runs $\{0.9,\;0.20187,\;-0.49627,\dots\}$. The test asserts a bit-identical map
across these values, with the offsets \emph{computed}, never hardcoded --- the spec's
rounded literals $(0.2019, -0.4963)$ carry $\sim7\times10^{-5}$ of residual $\alpha$,
which moves the phase at $k=30$ far above the assertion tolerance.
\end{observation}

\begin{remark}[The ambiguity is per-mode, and that is the right statement]
\label{rem:permode}
No single $\alpha' \ne \alpha$ aliases at \emph{every} $k$ simultaneously, since
\eqref{eq:aliasfamily} demands $\alpha'-\alpha = 2\pi n_k/(k^2\dt)$ with $n_k$ integer for
all $k$ at once, which fails already for two coprime $k$. So Theorem~\ref{thm:wrap} is
not a claim that the \emph{scalar} $\alpha$ is unidentifiable --- a model that has the
form $\omega = \alpha k^2 + \dots$ baked in can pin $\alpha$ down from low $k$ and then
predict high $k$ correctly. It is a claim about \textbf{nonparametric recovery of the
curve $\omega(k)$}: an unconstrained model has no reason to prefer one branch at large
$k$, because every branch fits the data exactly.

\emph{This is precisely why the theorem is interesting for a study about inductive bias,
and precisely why \code{K0} rather than \code{K2} is the headline rung.} \code{K2} is
handed the functional form and therefore escapes the ambiguity by assumption; \code{K0}
must discover it.
\end{remark}

\subsection{The positive theorem: the $\alpha$-derivative is wrap-free}
\label{sec:alphaderiv}

\begin{theorem}[$\alpha$-derivative identifiability]
\label{thm:alphaderiv}
For every $k$, including $k>\kwrap$,
\begin{equation}
\label{eq:alphaderiv}
\frac{\partial}{\partial\alpha}\arg m(k,\alpha) \;=\; -k^2\dt ,
\end{equation}
with $\beta$ and $V_0$ dropping out entirely. Moreover the derivative can be estimated
from \emph{wrapped} data without any unwrapping: for a grid $\alpha_0<\alpha_1<\dots$ with
gaps $\Delta\alpha_j$,
\[
\arg\bigl(m_{j+1}\,\overline{m_j}\bigr) \;=\; -\Delta\alpha_j\,k^2\dt
\qquad\text{exactly, provided}\qquad
\boxed{\;|\Delta\alpha_j|\;k^2\,\dt \;<\;\pi\;}
\]
\end{theorem}

\begin{proof}
From $\arg m(k,\alpha) = -(\alpha k^2 - \beta A^2 + V_0)\dt$ (as a real-valued lift),
differentiating in $\alpha$ gives \eqref{eq:alphaderiv}; the $\beta$ and $V_0$ terms are
$\alpha$-independent and vanish.

For the estimator: $m_{j+1}\overline{m_j}$ is a complex number whose true argument is
$-\Delta\alpha_j k^2\dt$. The function $\arg$ returns the representative in $(-\pi,\pi]$,
so it returns \emph{exactly} that value if and only if
$|\Delta\alpha_j k^2\dt| \le \pi$. Thus consecutive wrapped increments equal the true
increments under the stated condition, and no phase unwrapping --- and in particular no
reference to the truth --- is required.
\end{proof}

Theorem~\ref{thm:alphaderiv} is why identifiability survives above $\kwrap$ under
$\alpha$-varying supervision even though the \emph{absolute} phase does not. The
information is present in the data. Whether a given hypothesis class extracts it is an
empirical question --- experiment \textbf{G5b}, the central result of the study.

\begin{remark}[The precondition is checked, not assumed]
The estimator's precondition is verified on the caller's own $\alpha$ grid before any
measurement is taken, and the function raises otherwise. The dataset's realised gap is
$\Delta\alpha_{\max} = 3.467\times10^{-3}$; the requirement at $k=32$ is
$\Delta\alpha < \pi/(k^2\dt) = \pi/10.24 = 0.3068$. The margin is a factor of $\sim$88. A
\emph{probe} that chooses its own coarser spacing, however, can silently violate it and
return garbage --- hence the check.
\end{remark}

\subsection{$\alpha$-continuation is a different claim}
\label{sec:continuation}

Accumulating wrapped increments from an anchor gives the \emph{absolute} frequency:
\[
\arg m(k,\alpha_J) \;\overset{\text{lift}}{=}\; \arg m(k,\alpha_0)
\;+\;\sum_{j=0}^{J-1}\arg\bigl(m_{j+1}\overline{m_j}\bigr),
\qquad \omega = -\frac{(\cdot)}{\dt}.
\]

\begin{proposition}[Anchor condition]
\label{prop:anchor}
This recovers the true $\omega(k,\alpha_J)$ provided (i) the gap condition of
Theorem~\ref{thm:alphaderiv} holds on every step, and (ii) the anchor is itself unwrapped:
$|\omega(k,\alpha_0)|\,\dt \le \pi$.
\end{proposition}

Condition (ii) forces the anchor toward $\alpha_0\approx 0$ for large $k$ --- far outside
the training range $[0.7,1.1]$. As noted in \S\ref{sec:fno}, $\alpha = 0$ enters an FNO's
lift channel at $\tilde\alpha=-4.5$. \textbf{A failure of $\alpha$-continuation is
therefore evidence about $\alpha$-extrapolation, not about whether the
$\alpha$-derivative information was extracted}, and the two must be reported as separate
claims. Theorem~\ref{thm:alphaderiv} deliberately requires \emph{in-range} $\alpha$ only.

\subsection{Multi-$\dt$ as a second route past the horizon}

\begin{proposition}[Two step sizes]
\label{prop:multidt}
Observing the one-step multiplier at step sizes $\dt_1,\dt_2$ determines $\omega$ modulo
the generator of the lattice $\tfrac{2\pi}{\dt_1}\Z \,\cap\, \tfrac{2\pi}{\dt_2}\Z$,
i.e.\ modulo $\operatorname{lcm}\!\bigl(\tfrac{2\pi}{\dt_1},\tfrac{2\pi}{\dt_2}\bigr)$.
\end{proposition}

\begin{corollary}[A dyadic ladder buys only its smallest step]
\label{cor:dyadic}
For $\dt \in\{0.005,\,0.01,\,0.02\}$ the periods are $\{400\pi,\,200\pi,\,100\pi\}$, each
an integer multiple of the next, so their lcm is $400\pi$ --- the period of the smallest
$\dt$ alone. \textbf{The two coarser step sizes add no identifiability whatsoever.} The
horizon extends only as
\[
\kwrap(\dt) = \sqrt{\pi/(\alpha\dt)} \;\propto\; \dt^{-1/2},
\]
i.e.\ by $\sqrt2$ per halving: from $16.9$ at $\dt=0.01$ to $23.9$ at $\dt=0.005$
(for $\alpha=1.1$).
\end{corollary}

\begin{remark}
This is worth stating explicitly because it is counterintuitive and cheap to get wrong.
If the goal of the multi-$\dt$ arm (G6) were \emph{identifiability}, the step sizes should
be chosen \emph{incommensurate} --- then the lcm is enormous and the combined horizon far
exceeds either alone. As specified, G6's dyadic ladder is best understood as a test of
\textbf{$\dt$-transfer} (does a model trained at one $\dt$ work at another?), which is a
different and equally legitimate question; and $\dt$-transfer is meaningful only for the
split-step family, where $\dt$ genuinely multiplies a learned rate. An FNO ignores its
$\dt$ argument, so a transfer number measured on one is an artefact
(Remark~\ref{rem:fnodt}).
\end{remark}

\subsection{The horizon is arm-dependent}

$\kwrap$ is a function of $\alpha$, so widening the $\alpha$ range \emph{moves} the
horizon. On the production configuration $\alpha\in[0.7,1.1]$, $\dt=0.01$:
\[
\kwrap \in\bigl[\sqrt{\pi/0.011},\;\sqrt{\pi/0.007}\bigr] = [16.90,\,21.19].
\]
The G2 extrapolation arm ($\alpha\in[0.5,1.5]$) spans $[14.47,\,25.07]$. Every plot must
therefore draw \emph{its own arm's} horizon rather than a single global line, and the
banded-error edges must be recomputed per arm. The relevant scales for the production
configuration are
\[
\ktrain = 8 \;<\; \kwrap \in[16.9,\,21.2] \;<\; \knyq = 32 .
\]

\sourcenote{src/spno/evaluation/dispersion.py, src/spno/equations/nls.py, src/spno/data/shift.py}

\section{Spectral diagnostics}
\label{sec:spectraldiag}

Field-space $L^2$ hides both failures this study is about. A model can have small relative
$L^2$ while getting the \emph{dispersion relation} wrong --- error concentrated in a few
high-$k$ modes barely moves the norm --- and phase error is invisible to any metric that
looks only at $|\psi|$.

\subsection{Per-mode and banded error}

\[
E(k) \;=\; \Bigl\langle\, \bigl|\widehat{\psi}^{\mathrm{pred}}_k - \widehat{\psi}^{\mathrm{true}}_k\bigr| \,\big/\, \bigl|\widehat{\psi}^{\mathrm{true}}_k\bigr| \,\Bigr\rangle_{\text{batch}},
\qquad
E_{[a,b)} = \Bigl\langle \frac{\|\mathbf{1}_{[a,b)}(\widehat{\psi}^{\mathrm{pred}}-\widehat{\psi}^{\mathrm{true}})\|}{\|\mathbf{1}_{[a,b)}\widehat{\psi}^{\mathrm{true}}\|}\Bigr\rangle .
\]
Normalising per mode is what makes low-energy high-$k$ modes visible at all; the absolute
version is dominated by whichever modes carry the energy. Default band edges are
$\ktrain=8$ and $\kwrap=16.9$.

\subsection{Amplitude and phase separate exactly}

\begin{proposition}[Exact amplitude--phase decomposition]
\label{prop:ampphase}
With $p = \widehat\psi^{\mathrm{pred}}$, $t = \widehat\psi^{\mathrm{true}}$ and
$\Delta\varphi = \arg p - \arg t$,
\[
\underbrace{\sum_k|p_k-t_k|^2}_{\text{total}^2}
\;=\;
\underbrace{\sum_k\bigl(|p_k|-|t_k|\bigr)^2}_{\text{amplitude}^2}
\;+\;
\underbrace{2\sum_k |p_k||t_k|\bigl(1-\cos\Delta\varphi_k\bigr)}_{\text{phase}^2},
\]
and the phase term is nonnegative, vanishing exactly when the phases agree.
\end{proposition}

\begin{proof}
$|p-t|^2 = |p|^2+|t|^2-2\Real(p\bar t) = (|p|-|t|)^2 + 2|p||t| - 2|p||t|\cos\Delta\varphi$.
Sum over $k$.
\end{proof}

So taking the phase part as $\sqrt{\text{total}^2-\text{amplitude}^2}$, as
\code{amplitude\_and\_phase\_split} does, is not a heuristic residual. It is exactly
\[
\Bigl(2\textstyle\sum_k|p_k||t_k|\bigl(1-\cos\Delta\varphi_k\bigr)\Bigr)^{1/2}
\;\approx\;
\Bigl(\textstyle\sum_k|p_k||t_k|\,\Delta\varphi_k^2\Bigr)^{1/2}
\]
for small phase errors.

\begin{trap}[Never reduce a phase-error curve with \texttt{max}]
\label{trap:phasemax}
$\Delta\varphi_k$ is meaningful only where the target carries energy. In an \emph{empty}
mode the phase is roundoff noise, distributed uniformly on $(-\pi,\pi]$, so a bare
$\max_k$ reports $O(1)$ regardless of model quality. The scalar summary must be the
amplitude-weighted mean \emph{across modes},
\[
\overline{\Delta\varphi} \;=\; \frac{\sum_k |t_k|\,|\Delta\varphi_k|}{\sum_k|t_k|}.
\]
\end{trap}

\subsection{The G7 normalisation}

Fixing $\alpha$ makes the one-step task easier at \emph{every} $k$, so absolute $E(k)$
falls across the board and an absolute comparison cannot separate the three causes of
high-$k$ failure listed in \S\ref{sec:highk}. Only the ratio isolates cause (b):
\[
\text{band ratio} \;=\; \frac{E_{[\,\text{high},\infty)}}{E_{[0,\,\text{low})}},
\]
and even then only once cause (a) is removed by matching $m$ to Nyquist. This is a
\emph{numerical observation} on a trained model; it says nothing at untrained weights.

\subsection{The nonlinear cascade}

Initial conditions are band-limited to $|k|\le 8$ \emph{exactly} (a sharp cutoff, not a
Gaussian roll-off --- see \S\ref{sec:datadesign}), so every mode above the cutoff starts
empty and any energy found there is genuine nonlinear transfer through $\beta|\psi|^2\psi$.
A model that gets one-step $L^2$ right while failing to cascade is a failure that
field-space $L^2$ cannot see.

\begin{observation}[The cascade is real but stays below the horizon]
\label{obs:cascade}
On the training split: energy fraction above $\ktrain=8$ grows from $8.3\times10^{-6}$ at
step 1 to $2.7\times10^{-4}$ at step 200. Above $\kwrap$ it stays at
$\sim5\times10^{-7}$ to $2.4\times10^{-6}$ throughout. So in-distribution rollouts do
\emph{not} silently become spectral-extrapolation runs --- which is what makes G1 a clean
control. Measured, not assumed.
\end{observation}

\sourcenote{src/spno/evaluation/spectral.py, src/spno/data/generate.py}
\newpage
\part*{Part IV \quad The soft baseline and the misspecification dials}
\addcontentsline{toc}{part}{Part IV \; Soft baseline and misspecification}

\section{The PINO residual --- and the invariant it smuggles in}
\label{sec:pino}

\subsection{The discrete midpoint residual}

\begin{definition}[Crank--Nicolson / midpoint residual]
\label{def:cnresidual}
With $\overline\psi = \tfrac12(\psi_{n+1}+\psi_n)$ and
$\overline\rho = \tfrac12\bigl(|\psi_{n+1}|^2+|\psi_n|^2\bigr)$,
\[
R \;=\; \iu\,\frac{\psi_{n+1}-\psi_n}{\dt}
\;+\; \alpha\,\lap_h\overline\psi
\;+\; \beta\,\overline\rho\,\overline\psi
\;-\; V\,\overline\psi .
\]
The Laplacian is spectral, so it is exact for every mode the grid represents.
\end{definition}

The nonlinear term is averaged in the \emph{Delfour--Fortin--Payre} form
$\overline\rho\cdot\overline\psi$ rather than, say, $|\overline\psi|^2\overline\psi$. That
choice is what makes the scheme discretely mass-conserving --- and it is exactly the
problem.

\subsection{The confound, proved}

\begin{theorem}[Cayley transform: the CN update has modulus exactly one]
\label{thm:cayley}
Let $\psi_n = A\eu^{\iu kx}$ with constant $V\equiv V_0$, and let $\psi_{n+1}=z\psi_n$ with
$|z|=1$ (so that $\overline\rho = A^2$). Then $R\equiv 0$ if and only if
\begin{equation}
\label{eq:cayley}
z \;=\; \frac{1-\iu\,\omega\dt/2}{1+\iu\,\omega\dt/2},
\qquad \omega = \alpha k^2 - \beta A^2 + V_0,
\end{equation}
and this $z$ satisfies $|z| = 1$ \textbf{exactly}, for every $\dt$ and every $\omega$.
\end{theorem}

\begin{proof}
With $\psi_{n+1}=z\psi_n$, $\overline\psi = \tfrac{1+z}{2}\psi_n$, $\lap_h\psi_n = -k^2\psi_n$
and $\overline\rho = A^2$:
\[
\frac{R}{\psi_n} \;=\; \iu\,\frac{z-1}{\dt} \;+\; \bigl(-\alpha k^2 + \beta A^2 - V_0\bigr)\frac{1+z}{2}
\;=\; \iu\,\frac{z-1}{\dt} \;-\; \omega\,\frac{1+z}{2}.
\]
Setting this to zero and writing $a=\omega\dt/2$ gives $\iu(z-1) = a(1+z)$, hence
$z(\iu - a) = a+\iu$ and
\[
z \;=\; \frac{a+\iu}{\iu-a} \;=\; \frac{1-\iu a}{1+\iu a},
\]
which is \eqref{eq:cayley}. Its modulus is $|1-\iu a|/|1+\iu a| = \sqrt{1+a^2}/\sqrt{1+a^2}=1$.
\end{proof}

\begin{corollary}[The physics loss contains the invariant under study]
\label{cor:pinoconfound}
On plane waves the CN residual is exactly solved by a \emph{unitary} update. More
generally the DFP-averaged scheme conserves the discrete mass. \textbf{The 7a physics
loss therefore smuggles in the very invariant the study compares methods on.}
\end{corollary}

This is not a reason to avoid the baseline --- it is precisely what the
loss/projection/architecture distinction exists to expose --- but it must be stated in the
thesis \emph{next to every 7a number}.

\begin{corollary}[The scheme's own frequency]
\label{cor:cnfreq}
Writing $z = \eu^{-\iu\omega_{\mathrm{CN}}\dt}$ and inverting \eqref{eq:cayley},
\[
\boxed{\;\omega_{\mathrm{CN}} \;=\; \frac{2}{\dt}\arctan\!\Bigl(\frac{\omega\dt}{2}\Bigr)\;}
\]
\end{corollary}

\begin{proof}
$\arg\frac{1-\iu a}{1+\iu a} = -2\arctan a$ with $a=\omega\dt/2$, so
$\omega_{\mathrm{CN}}\dt = 2\arctan(\omega\dt/2)$.
\end{proof}

\begin{trap}[A test that feeds the analytic plane wave and expects zero tests the wrong
thing]
The scheme's own exact solution advances at $\omega_{\mathrm{CN}}$, not at $\omega$. The
gap between them \emph{is} the scheme's $O(\dt^2)$ truncation error, and a residual test
must be built around Corollary~\ref{cor:cnfreq}, not around $\eu^{-\iu\omega\dt}$.
\end{trap}

\subsection{The residual is second order --- with the coefficient derived}

\begin{proposition}[Residual expansion]
\label{prop:residualexpansion}
Feeding the residual the \emph{analytic} pair $\psi_{n+1} = \eu^{-\iu\omega\dt}\psi_n$
gives, with $\theta = \omega\dt$,
\[
\frac{R}{\psi_n}
\;=\; \frac{1}{\dt}\Bigl[\tfrac{1}{12}\theta^3 \;-\; \tfrac{\iu}{24}\theta^4 \;-\; \tfrac{1}{80}\theta^5 \;+\; O(\theta^6)\Bigr],
\]
that is
\[
\Real\!\Bigl(\frac{R}{\psi_n}\Bigr) = \frac{\omega^3\dt^2}{12} + O(\dt^4),
\qquad
\Imag\!\Bigl(\frac{R}{\psi_n}\Bigr) = -\frac{\omega^4\dt^3}{24} + O(\dt^5).
\]
The next real term is $-\omega^5\dt^4/80$, so the leading coefficient is recovered only up
to a \emph{relative} correction of order $\tfrac{12}{80}(\omega\dt)^2 = 0.15\,(\omega\dt)^2$.
\end{proposition}

\begin{proof}
From the proof of Theorem~\ref{thm:cayley} with $z=\eu^{-\iu\theta}$,
\[
\frac{R}{\psi_n}\;=\;\frac{1}{\dt}\Bigl[\iu\bigl(\eu^{-\iu\theta}-1\bigr) \;-\; \frac{\theta}{2}\bigl(1+\eu^{-\iu\theta}\bigr)\Bigr].
\]
Expand $\eu^{-\iu\theta} = 1 - \iu\theta - \tfrac{\theta^2}{2} + \tfrac{\iu\theta^3}{6} + \tfrac{\theta^4}{24} - \tfrac{\iu\theta^5}{120}+\cdots$. Then
\[
\iu\bigl(\eu^{-\iu\theta}-1\bigr) = \theta - \tfrac{\iu\theta^2}{2} - \tfrac{\theta^3}{6} + \tfrac{\iu\theta^4}{24} + \tfrac{\theta^5}{120},
\]
\[
\tfrac{\theta}{2}\bigl(1+\eu^{-\iu\theta}\bigr) = \theta - \tfrac{\iu\theta^2}{2} - \tfrac{\theta^3}{4} + \tfrac{\iu\theta^4}{12} + \tfrac{\theta^5}{48}.
\]
Subtracting, the $\theta$ and $\theta^2$ terms cancel identically and
\[
\theta^3\Bigl(-\tfrac16+\tfrac14\Bigr) + \iu\theta^4\Bigl(\tfrac1{24}-\tfrac1{12}\Bigr) + \theta^5\Bigl(\tfrac1{120}-\tfrac1{48}\Bigr)
= \tfrac{\theta^3}{12} - \tfrac{\iu\theta^4}{24} - \tfrac{\theta^5}{80}.
\]
Dividing by $\dt$ and substituting $\theta = \omega\dt$ gives the stated expansion.
\end{proof}

\subsection{Normalising the loss}

\[
\mathcal{L}_{\mathrm{phys}} \;=\; \Bigl\langle \frac{\|R\|_{L^2_h}}{\|\psi_n\|_{L^2_h}\,/\,\dt}\Bigr\rangle_{\text{batch}} .
\]
The normalisation makes the weight $\lambda$ comparable across $\dt$ and across mass ---
both of which vary by design. An \emph{un}normalised residual weight becomes
uninterpretable the moment the multi-$\dt$ arm exists, and would quietly turn the
objective into a mass-weighted one. The $\lambda$ sweep $\{0.01,0.1,1,10\}$ is reported in
full, never as a tuned value.

\sourcenote{src/spno/losses/pde\_residual.py}

\section{Misspecification --- moving the truth outside the constrained class}
\label{sec:misspec}

Phase 9 is the only phase that makes ``structure beats FNO'' a \emph{question} rather than
a tautology, because it moves the data-generating equation outside the hypothesis class of
the structured models. Two dials, each breaking a different assumption, each recovering
exact NLS at zero.

\subsection{Dial 1 --- nonlocal nonlinearity (breaks locality)}

\begin{definition}[Nonlocal generator]
\label{def:nonlocal}
Replace the local rate by
\[
\nu \;=\; \beta\,(W_\sigma * \rho) \;-\; V,
\qquad
\widehat{W_\sigma}(k) \;=\; \eu^{-\sigma^2|k|^2/2}.
\]
\end{definition}

\begin{proposition}[The kernel is a unit-mass symmetric mollifier for every $\sigma$]
\label{prop:kernel}
$\widehat{W_\sigma}(0) = 1$ exactly, so $\int W_\sigma = 1$ and the kernel never rescales
the mean density; $\widehat{W_\sigma}$ is real and even in $k$, so $W_\sigma$ is real and
symmetric; and $\sigma\to0$ gives $\widehat{W}\to1$, i.e.\ a delta, \emph{by construction}
rather than in a limit that has to be argued.
\end{proposition}

\begin{guarantee}[Mass and $U(1)$ survive]
$\nu$ is real, so the local substep is still a modulus-one multiplier: discrete mass is
preserved exactly and global $U(1)$ equivariance survives, for every $\sigma$. Model B's
projection therefore remains \emph{correct} under this dial.
\end{guarantee}

\begin{theorem}[The dynamics remain Hamiltonian]
\label{thm:nonlocalham}
For symmetric $W$, the local substep of Definition~\ref{def:nonlocal} is still the exact
flow of a sub-Hamiltonian, namely
\[
\Ham_N \;=\; -\frac{\beta}{2}\iint \rho(x)\,W(x-y)\,\rho(y)\dd x\dd y \;+\; \int V\rho\dd x .
\]
Hence the composition is still symplectic, and $\rho$ is still frozen along the flow.
\textbf{What is broken is locality, and nothing else.}
\end{theorem}

\begin{proof}
Let $G[\rho] = \tfrac{\beta}{2}\iint\rho(x)W(x-y)\rho(y)$. Then
\[
\frac{\delta G}{\delta\rho(x)} = \frac{\beta}{2}\Bigl(\int W(x-y)\rho(y)\dd y + \int \rho(y)W(y-x)\dd y\Bigr)
= \beta\,(W*\rho)(x),
\]
where the two terms combine into one \emph{only because} $W(x-y)=W(y-x)$. So
$\delta(-G+\int V\rho)/\delta\rho = -\bigl[\beta(W*\rho) - V\bigr] = -\nu$, and by the
computation in the proof of Theorem~\ref{thm:learnedmain}(3),
$\delta\Ham_N/\delta\bar\psi = (\delta\Ham_N/\delta\rho)\,\psi = -\nu\psi$, which generates
$\partial_t\psi = \iu\nu\psi$. Since $\nu$ is real, $\rho$ is frozen
(Lemma~\ref{lem:local}) and the closed-form multiplier is the exact flow.
\end{proof}

\begin{remark}[Why this dial separates C1 from C2]
C1's $\nu_\theta$ is \emph{pointwise} in $\rho$ and therefore \textbf{cannot represent}
$\beta(W_\sigma*\rho)$ for $\sigma>0$: the truth has left C1's class. C2's $\nu_\theta$ is
an FNO on $\rho$ and \emph{can}. This is the cleanest available separation of the two, and
it is why the dial was chosen.
\end{remark}

\begin{remark}[Deliberately \emph{not} a saturable or quintic nonlinearity]
$\nu_\theta$ is already a free function of the scalar $\rho$, so the C family learns
$\nu=\beta\rho/(1+s\rho)$ or $\nu = \beta\rho + c\rho^2$ easily. Those are not
misspecifications at all --- they are inside the class. Only breaking \emph{locality} or
breaking \emph{conservation} moves the truth out.
\end{remark}

\subsection{Dial 2 --- weak gain/loss (breaks conservation itself)}

\begin{definition}[Gain/loss generator]
\label{def:gainloss}
\[
\iu\psi_t + \alpha\lap\psi + \beta|\psi|^2\psi - V\psi \;=\; +\,\iu\gamma\,\psi ,
\qquad\text{i.e.}\qquad
\psi_t = \gamma\psi + \iu\bigl(\alpha\lap\psi+\beta|\psi|^2\psi-V\psi\bigr).
\]
\end{definition}

\begin{theorem}[Mass grows at a fixed exponential rate]
\label{thm:gainloss}
Along \eqref{def:gainloss},
\[
\boxed{\;\frac{\dd}{\dd t}\ln\Mass \;=\; 2\gamma\;}
\]
\end{theorem}

\begin{proof}
$\frac{\dd}{\dd t}\!\int|\psi|^2 = \int 2\Real(\bar\psi\psi_t)
= \int 2\Real\bigl(\gamma|\psi|^2\bigr) + \int 2\Real\bigl(\iu\bar\psi(\cdots)\bigr)$.
The second integral vanishes by the computation in Theorem~\ref{thm:mass}. The first is
$2\gamma\Mass$. Hence $\dd\Mass/\dd t = 2\gamma\Mass$.
\end{proof}

\begin{remark}[The sign convention is pinned by a measured rate, not by prose]
Writing $-\iu\gamma\psi$ on the right-hand side would give decay instead. The tests assert
$\dd\ln\Mass/\dd t = 2\gamma$ as a \emph{rate} rather than trusting the derivation ---
which is the correct level of paranoia for a sign that is this easy to invert.
\end{remark}

\begin{remark}[This dial makes Model B actively wrong]
Under $\gamma\neq0$ the true mass is \emph{not} conserved, so Model B's hard projection is
no longer merely unhelpful --- it enforces a false constraint. A located crossover here is
direct evidence about \emph{when a hard invariant stops helping}, which is the
generalisation the whole study is aiming at. A \emph{bounded} result (``no crossover
within the swept range'') is equally publishable.
\end{remark}

\begin{proposition}[What is exactly, and what is only approximately, exact]
\label{prop:gainexact}
The gain factor $\eu^{\gamma\tau}$ is a \emph{real scalar} and therefore commutes with the
linear Fourier multiplier $\Kop_\tau$. Consequently splitting the gain away from the
kinetic step introduces \textbf{no additional splitting error}, and the dial does not
contaminate $\epssplit$.

It does not, however, commute with the nonlinear local step. Along the exact local
sub-flow $\psi_t = \gamma\psi + \iu(\beta\rho-V)\psi$ one has $\rho(\tau) = \rho_0\eu^{2\gamma\tau}$,
so the exact accumulated phase is
\[
\int_0^{\tau}\!\bigl(\beta\rho_0\eu^{2\gamma s} - V\bigr)\dd s
= \beta\rho_0\,\tau\bigl(1+\gamma\tau\bigr) - V\tau + O(\gamma^2\tau^3),
\]
whereas freezing $\rho$ at $\rho_0$ gives $\beta\rho_0\tau - V\tau$. The local substep
therefore carries a phase error $\beta\rho_0\gamma\tau^2 + O(\gamma^2\tau^3)$ --- second
order in the substep and \emph{first order in the dial}, hence vanishing as $\gamma\to0$
and further suppressed by a factor $1/M$ by the substepping.
\end{proposition}

The precise statement to make in the thesis is therefore: \emph{gain and kinetic commute
exactly; the residual $O(\beta\gamma\dt^2/M)$ term lives inside the local substep and is
suppressed by the substepped reference along with the ordinary Strang commutator.}

\subsection{The dial-zero gate}

\begin{trap}[A spectral Gaussian at $\sigma=0$ is not the identity]
\label{trap:dialzero}
$\widehat{W_0}\equiv 1$ exactly, but $\mathcal{F}^{-1}(\mathcal{F}\rho \cdot 1)$ is
\textbf{not bitwise} $\rho$ --- the FFT round trip is accurate, not exact. Phase 9's
opening gate asserts that the dial-zero generator reproduces the unperturbed solver
\emph{bitwise}, so both perturbed operators short-circuit to the plain Strang forward pass
at zero, and the exact arm is routed through the \emph{same class} the production shards
were generated with. This is load-bearing, not an optimisation: a crossover measured
against a subtly different zero point is uninterpretable.
\end{trap}

\begin{trap}[Turn one dial at a time]
Two simultaneous perturbations cannot be attributed to either broken assumption. The
configuration object refuses the combination outright.
\end{trap}

\begin{remark}[The perturbed generators must be substepped too]
Both dials are wrapped in the same $M=32$ substepping as the unperturbed reference. A
single perturbed Strang step at $\dt$ would confound the dial with the $O(\dt^2)$
splitting error, and every located crossover would then be partly a measurement of
$\epssplit$.
\end{remark}

\begin{remark}[Do not add other equations]
NLS was chosen because it bundles nonlinearity, complex fields, spectral structure,
Hamiltonian invariants, phase-sensitive error and long-horizon stability into one
equation. Adding Burgers, wave or Navier--Stokes before Phase 9 is complete costs the
thesis its spine.
\end{remark}

\sourcenote{src/spno/solvers/perturbed.py, src/spno/misspecification.py}
\newpage
\part*{Part V \quad Data, measurement, and the phase map}
\addcontentsline{toc}{part}{Part V \; Data, measurement, phase map}

\section{The training distribution, and why it is shaped this way}
\label{sec:datadesign}

Two design choices exist purely to keep later experiments honest.

\begin{enumerate}[leftmargin=1.4em,itemsep=0.4em]
\item \textbf{Sharp spectral cutoff.} Initial conditions are band-limited \emph{exactly}
to $|k|\le\ktrain=8$:
\[
\widehat{\psi_0}(k) \;=\; \mathbf{1}_{|k|\le \ktrain}\;\eu^{-|k|^2/(2\ktrain^2)}\;\widehat{\xi}(k),
\qquad \xi \text{ white noise},
\]
with real and imaginary parts drawn independently (the NLS state is a genuine complex
field, not a pair of coupled real ones). A \emph{soft} filter that leaked high-$k$ energy
into the training set would void the G4 bandwidth-shift design with no visible symptom.

\item \textbf{Mass varies across samples} --- drawn uniformly in $[1,3]$, a $3\times$
range. Normalising every initial condition to unit mass would make mass conservation
\emph{memorisable}, which flatters the unconstrained baselines and turns Model B's
projection into a no-op.
\end{enumerate}

Potentials are smooth, real, and normalised so that $\max|V|$ equals the sampled
amplitude, with amplitude, correlation length and functional family controlled
\emph{independently} so that the G3 potential arm can vary one at a time.

\subsection{The resolution-adequacy screen}

Focusing NLS ($\beta>0$) can steepen a field until the grid stops resolving it, and the
resulting error looks exactly like an out-of-distribution failure. The screen measures the
spectral tail fraction above $0.75\,\knyq$ and requires it to stay below $10^{-6}$;
trajectories that fail are excluded, or the $(\alpha,\beta)$ box is narrowed.

\begin{observation}[The production box is resolved with enormous margin]
Maximum tail fraction over 64 screened samples: $2.52\times10^{-10}$, median
$3.55\times10^{-13}$, against a threshold of $10^{-6}$. Zero samples over threshold.
Verdict: \emph{resolved}.
\end{observation}

\section{Resolution transfer, honestly}
\label{sec:resolution}

\subsection{Spectral resampling}

\begin{proposition}[Exact resampling of band-limited fields]
\label{prop:resample}
Moving a periodic field between grids by truncating or zero-padding its spectrum is
\textbf{exact} for any field band-limited below both Nyquists. With the
\code{norm="backward"} convention the inverse transform carries the $1/N$, so coefficients
must be rescaled by $N_{\mathrm{tgt}}/N_{\mathrm{src}}$ to hold the \emph{physical}
amplitude fixed rather than the discrete one.
\end{proposition}

\begin{proof}
$\psi(x_j) = \tfrac{1}{N}\sum_k\widehat\psi_k\eu^{\iu kx_j}$. For the continuum function to
be unchanged on the target grid one needs
$\tfrac{1}{N'}\sum_k \widehat\psi'_k\eu^{\iu kx} = \tfrac{1}{N}\sum_k\widehat\psi_k\eu^{\iu kx}$
for all $x$, hence $\widehat\psi'_k = (N'/N)\widehat\psi_k$ on the common modes and zero
elsewhere. Exactness follows because a band-limited field has no coefficients outside that
common set.
\end{proof}

\begin{trap}[The Nyquist mode is self-conjugate, and cannot be upsampled]
\label{trap:nyquist}
On an even grid the mode $k=N/2$ is self-conjugate: it represents $\cos(Nx/2)$ with no way
to distinguish $+k$ from $-k$. Splitting it between the two target modes on a finer grid
is a \textbf{choice, not a fact}, and silently making one would inject an arbitrary phase
into every downstream number. Upsampling a field with relative energy above $10^{-12}$ at
the source Nyquist is therefore refused rather than resolved by convention.
\end{trap}

\subsection{Which parameters are grid-independent}

\begin{center}
\begin{tabular}{@{}p{0.27\linewidth}p{0.66\linewidth}@{}}
\toprule
\textbf{parameter group} & \textbf{status}\\
\midrule
FNO \code{spectral.weight} & grid-independent \emph{for $|k|\le m$}: indexed by frequency,
and \code{irfft} receives $n$ at call time. This is the \textbf{only} source of the FNO's
resolution transfer, and it says nothing about modes above the budget.\\
FNO lift / local / project & grid-independent: pointwise in the channel axis.\\
C-family local phase & grid-independent: pointwise in $(\rho,V,\alpha,\beta)$, no
grid-shaped parameter. (C2's phase FNO likewise indexes by frequency.)\\
C-family \code{kinetic.k\_squared} & \emph{grid-dependent input range} --- see
Proposition~\ref{prop:frozennorm}.\\
Mass projection & grid-independent: rescaling by a mass ratio involves no grid-shaped
parameter.\\
\bottomrule
\end{tabular}
\end{center}

\begin{proposition}[The frozen normaliser, and its honest cost]
\label{prop:frozennorm}
The kinetic MLP consumes $\widetilde{k^2} = |k|^2/\max_{\mathrm{grid}}|k|^2$. On rebinding
to a finer grid the normaliser is \textbf{frozen at the training grid's value}. Freezing is
necessary: recomputing it would send the same \emph{physical} $k$ to a different network
input ($1024$ at $N=64$, $4096$ at $N=128$), silently changing $\kappa$ everywhere without
touching a weight. The cost, which must be reported rather than hidden, is that modes above
the training grid's Nyquist then enter at $\widetilde{k^2}>1$ --- outside the input range
the MLP ever saw.
\end{proposition}

\begin{remark}[The genuine resolution test is G4 wearing a different hat]
Upsampling a \emph{band-limited} field and re-running an FNO barely changes its output,
because the FNO only touches $|k|\le m$ and those modes are unchanged by the resample.
\textbf{That is not evidence of resolution transfer.} The genuine test is new energy above
the training grid's Nyquist --- which is spectral extrapolation (G4), not resolution
transfer. The two arms are reported separately, with this sentence attached.
\end{remark}

\sourcenote{src/spno/evaluation/resolution.py}

\section{Measurement: precision, drift classification, and order estimation}
\label{sec:measurement}

\subsection{Precision discipline}

Training runs in \code{float32}/\code{complex64} for speed. \textbf{Anything measuring a
conservation law is explicitly widened to \code{float64}/\code{complex128} and evaluated
on CPU}, so that arithmetic roundoff never masks --- or manufactures --- a genuine
architectural violation.

\begin{observation}[The two floors, quoted beside every ``exact'' claim]
\label{obs:floors}
Mass drift of the reference over 100 steps:
\[
\text{\code{float64}, CPU}: 1.87\times10^{-14},
\qquad
\text{\code{float32}, CPU}: 1.43\times10^{-5}.
\]
The \code{float32} drift accumulates \textbf{linearly}, at $\approx1.2\times10^{-7}$ per
split step --- not as a random walk. That is the signature of a systematic (biased)
round-off in the FFT round trip rather than a stochastic one, and it means the drift is
predictable from the step count. MPS is $\approx3\times$ worse than CPU.
\end{observation}

\begin{trap}[Floors must be measured with one split step per model step]
The trajectory generator takes $M=32$ substeps per stored step. Measuring a
\code{float32} mass floor with the generator's substepping would report $32\times$ the
drift a model actually incurs. Use one split step per model step.
\end{trap}

\begin{trap}[Three ways to destroy a complex network while widening it]
\label{trap:widen}
(i) \code{.double()} on a complex tensor \textbf{silently drops the imaginary part}.
(ii) \code{.double()} on a module \textbf{skips} complex parameters entirely.
(iii) \code{.to(torch.float64)} \textbf{casts the complex spectral weights to real},
destroying the operator with only a warning. All three are converted into hard exceptions
by an explicit dtype guard inside the spectral convolution, and widening goes through a
single dedicated helper.
\end{trap}

\begin{trap}[\texttt{model.to(device)} mutates in place]
It will strand a core that another model shares --- and Model B-post deliberately reuses
Model A's weights.
\end{trap}

\subsection{Bounded versus secular: classification by fitted slope}
\label{sec:drift}

``Bounded'' and ``secular'' differ by whether drift keeps \emph{growing with the horizon},
which is a statement about a trend, so it is fitted rather than eyeballed:
\[
\log d_i \;\approx\; c + s\,\log n_i,
\qquad
\text{classification} = \begin{cases}\text{bounded} & s < 0.25\\ \text{secular} & s\ge 0.25\end{cases}
\]
with $s$ reported together with its standard error. A conserved-up-to-model-error quantity
plateaus ($s\approx0$); linear accumulation gives $s\approx1$. The cut at $0.25$ is
deliberately \textbf{generous to ``bounded''}, so that a borderline case is reported as
secular rather than flattering the structured model. Points with $d_i\le0$ are dropped:
they are exact conservation, which has no slope.

\subsection{Order estimation}

Every order claim in the project is a measured $\log_2$ ratio compared against the
expected value, never a threshold:
\[
p_i = \log_2\frac{e_{i+1}}{e_i},
\qquad
\text{regime} = \begin{cases}
\text{exact} & \max_i e_i < 10^{-12}\ \text{and flat in }\dt\\
\text{approximate, order } \bar p & \max_i p_i - \min_i p_i < 0.25\\
\text{inconsistent} & \text{otherwise}
\end{cases}
\]
Varying $\dt$ away from the trained value is deliberate here: it probes the
\emph{architecture}, and is not a claim about prediction accuracy at those step sizes.
Because a trained model legitimately refuses off-$\dt$ evaluation
(Remark~\ref{rem:fnodt}), the licence to do so is a scoped, temporary opt-in rather than a
flag that stays flipped.

\subsection{Run-identifier hygiene}

Every swept knob enters the run identifier. A four-value $\lambda$ sweep writing to one
directory silently overwrites three results. The pattern is
\code{f"\{config\_hash(data\_config)\}-\{arm\}-\{knob\}"}, and the misspecification
identifier deliberately \emph{collapses} to the bare data hash at dial zero so the exact
arm reuses the existing shards untouched while every perturbed arm gets its own directory.

Gates use \code{raise RuntimeError}, never \code{assert}: \code{python -O} strips asserts,
and a gate is the only thing standing between an instrument bug and a full run of
plausible-but-fictional numbers.

\section{The phase map}
\label{sec:phasemap}

Which mathematics each phase of the pipeline rests on.

\begin{center}
\small
\begin{longtable}{@{}p{0.055\linewidth}p{0.30\linewidth}p{0.58\linewidth}@{}}
\toprule
\textbf{Ph.} & \textbf{What it does} & \textbf{Mathematics it rests on}\\
\midrule
\endhead
\textbf{0} & Reference validation & Lemmas~\ref{lem:kinetic}--\ref{lem:local} (exact
sub-flows); Theorems~\ref{thm:strangmass}--\ref{thm:strangorder} (mass, reversibility,
symplecticity, order 2); Corollary~\ref{cor:modham} (bounded energy);
Definition~\ref{def:epssplit} and \S\ref{sec:fittedfloor} (the floor and the fitted floor);
\S\ref{sec:datadesign} (resolution screen).\\[0.4em]

\textbf{1} & Data generation & \S\ref{sec:datadesign} (sharp cutoff, varying mass, potential
families); Remark~\ref{rem:tautology} (targets must be substepped);
Observation~\ref{obs:cascade} (the cascade stays below $\kwrap$);
Theorem~\ref{thm:alphaderiv} precondition ($\Delta\alpha\,k^2\dt<\pi$) checked on the
realised draw.\\[0.4em]

\textbf{2--3} & Baselines A and B & \S\ref{sec:fno} (FNO structure, grid independence,
the three high-$k$ causes); Proposition~\ref{prop:metricproj} (projection is the metric
projection); Guarantee~\ref{guar:B}; \S\ref{sec:drift} (secular vs bounded classification).\\[0.4em]

\textbf{4--5} & Structured models C1/C2/C3 & Theorem~\ref{thm:learnedmain} (the four
guarantees, with proofs); Corollary~\ref{cor:boundedenergy} (the falsifiable prediction);
Guarantees~\ref{guar:c2}--\ref{guar:c3} and the C3 order-2 proof (\S\ref{sec:c3});
Remark~\ref{rem:kam} (reversible-KAM caveat); \S\ref{sec:ladders} (the K/L ladders and
Observation~\ref{obs:ladder}).\\[0.4em]

\textbf{6} & Generalisation and identifiability \emph{(the spine)} &
Definition~\ref{def:multiplier} and Proposition~\ref{prop:strangexactpw} (the probe, and
why its floor is roundoff not $\epssplit$); \textbf{Theorem~\ref{thm:wrap}} (G5a: the
negative theorem) with Corollary~\ref{cor:multistep} and
Proposition~\ref{prop:alias}/Remark~\ref{rem:permode};
\textbf{Theorem~\ref{thm:alphaderiv}} (G5b: the positive theorem);
Proposition~\ref{prop:anchor} (why continuation is a \emph{different} claim);
Corollary~\ref{cor:dyadic} (G6: what a dyadic $\dt$ ladder actually buys);
\S\ref{sec:probeamp} (G-arms: probe amplitude and the mass confound);
\S\ref{sec:spectraldiag} (G7 band ratio; G9 cascade).\\[0.4em]

\textbf{7} & PINO baseline & Definition~\ref{def:cnresidual};
\textbf{Theorem~\ref{thm:cayley}} and Corollary~\ref{cor:pinoconfound} (the physics loss
contains the invariant under study); Corollary~\ref{cor:cnfreq} (the scheme's own
frequency); Proposition~\ref{prop:residualexpansion} (second order, coefficients derived);
loss normalisation.\\[0.4em]

\textbf{8} & Resolution and robustness & Proposition~\ref{prop:resample} and
Trap~\ref{trap:nyquist} (exact resampling; the Nyquist obstruction);
Proposition~\ref{prop:frozennorm} (the frozen normaliser and its cost);
Proposition~\ref{prop:fnogrid} (what grid independence does and does not mean);
\S\ref{sec:measurement} (noise, sparsity, and the precision floors).\\[0.4em]

\textbf{9} & Misspecification sweep & Definition~\ref{def:nonlocal},
Proposition~\ref{prop:kernel}, \textbf{Theorem~\ref{thm:nonlocalham}} (nonlocal but still
Hamiltonian --- what separates C1 from C2); Definition~\ref{def:gainloss} and
\textbf{Theorem~\ref{thm:gainloss}} ($\dd\ln\Mass/\dd t = 2\gamma$);
Proposition~\ref{prop:gainexact} (what is exactly vs approximately exact under the dial);
Trap~\ref{trap:dialzero} (the bitwise gate).\\
\bottomrule
\end{longtable}
\end{center}

\section{Constants and measured values}
\label{sec:constants}

All at configuration hash \code{bd4e108527}: $N=64$, $\dt=0.01$, $M=32$ substeps,
$\alpha\sim U(0.7,1.1)$, $\beta\sim U(-0.4,0.6)$, $\psi_0$ band-limited to $|k|\le8$,
mass $\sim U(1,3)$, $800/100/100$ disjoint trajectories $\times$ $200$ steps
$=$ $160{,}000$ one-step training pairs.

\begin{center}
\begin{tabular}{@{}p{0.52\linewidth}r@{}}
\toprule
\textbf{quantity} & \textbf{value}\\
\midrule
$\epssplit$ (splitting floor, mean) & $6.239\times10^{-5}$\\
\hspace{1em} median / min / max & $4.484\times10^{-5}$ / $1.155\times10^{-6}$ / $2.801\times10^{-4}$\\
$\epssplit^{\mathrm{fit}}/\epssplit$, scalar family & $0.9999978$\\
$\epssplit^{\mathrm{fit}}/\epssplit$, per-mode (\code{K2}) family & $0.99999$\\
fitted kinetic / local scale & $1.0000008$ / $0.9999979$\\
Strang convergence order (finest ratio) & $2.0465$\\
reference self-convergence, $M{=}32$ vs $M{=}64$ & $1.50\times10^{-7}$ ($0.24\%$ of $\epssplit$)\\
\midrule
$\ktrain$ / $\kwrap$ / $\knyq$ & $8$ / $16.90$--$21.19$ / $32$\\
$\alpha$ max gap (train) & $3.467\times10^{-3}$ \; (needs ${<}\,0.3068$)\\
probe amplitude $A$ & $0.5642$\\
\midrule
mass floor, \code{float64} CPU (100 steps) & $1.87\times10^{-14}$\\
mass floor, \code{float32} CPU (100 steps) & $1.43\times10^{-5}$\\
\code{float32} drift per split step & $\approx1.2\times10^{-7}$, \emph{linear}\\
long run: max mass drift, \code{float64} (1000 steps) & $3.82\times10^{-13}$\\
long run: max energy drift; classification & $7.92\times10^{-5}$; \emph{bounded}\\
resolution screen: max tail fraction & $2.52\times10^{-10}$ (threshold $10^{-6}$)\\
\bottomrule
\end{tabular}
\end{center}

\subsection*{Phases 2--3 (3 seeds, 25 epochs, $287{,}746$ params each)}

\begin{center}
\begin{tabular}{@{}lrrrrl@{}}
\toprule
\textbf{model} & \textbf{1-step} & \textbf{roll@100} & \textbf{mass@100} & \textbf{energy@100} & \textbf{trend}\\
\midrule
A \; (FNO) & $6.94\times10^{-4}$ & $2.77\times10^{-2}$ & $1.61\times10^{-2}$ & $1.79\times10^{-2}$ & secular\\
B-post & $6.74\times10^{-4}$ & $2.57\times10^{-2}$ & $1.3\times10^{-15}$ & $9.19\times10^{-3}$ & secular\\
B-loop & $6.62\times10^{-4}$ & $2.44\times10^{-2}$ & $1.2\times10^{-15}$ & $8.14\times10^{-3}$ & secular\\
\bottomrule
\end{tabular}
\end{center}

The FNO does \textbf{not} beat $\epssplit$; it sits $11\times$ above it, so the structured
models have headroom.

\subsection*{Phases 4--5 (guarantees verified at \emph{random untrained} weights)}

\begin{center}
\begin{tabular}{@{}lrrrl@{}}
\toprule
\textbf{model} & \textbf{params} & \textbf{vs A} & \textbf{mass @init} & \textbf{reversibility}\\
\midrule
C1 (pointwise $\rho$ phase) & $2{,}466$ & $0.009\times$ & $4.4\times10^{-16}$ & exact ($\dt$-independent)\\
C2 (FNO-on-$\rho$ phase) & $288{,}770$ & $1.004\times$ & $2.2\times10^{-16}$ & exact\\
C3 (Re/Im $\psi$ phase) & $288{,}834$ & $1.004\times$ & $4.4\times10^{-16}$ & \textbf{order 2.00}\\
\bottomrule
\end{tabular}
\end{center}

\begin{center}
\emph{A vs C2 is the fair fight: same backbone, matched capacity, differing only in whether
the network output is the field or a phase applied to it.}
\end{center}

\vspace{1.5em}
\hrule
\vspace{0.8em}

\subsection*{A note on outstanding debt}

Phases 2--3 as tabulated were \textbf{budget-bound}: \code{best\_epoch == 24} for all six
runs (early stopping never fired) with validation loss still falling $26$--$50\%$ over the
last five epochs. The numbers above are therefore a lower bound on all three models'
achievable accuracy and are not yet convergence-comparable. Nothing in this document's
\emph{mathematics} depends on them; the guarantees are architectural and hold at random
initialisation. But no accuracy number should enter the thesis until the budget-bound
warning is empty for every model in the comparison.

\end{document}
